% =====================================================================
%  CAPITULO 4 - DIVISOR DE TENSAO E CORRENTE
%  (versao revisada e ampliada: 2 -> 16 questoes)
% =====================================================================
\section{Divisor de Tensão e de Corrente}

\begin{conceito}[Antes de começar]
\textbf{Divisor de tensão} (resistores em \emph{série}): a tensão total se
reparte na proporção direta das resistências.
\[
U_k = U_{\text{total}}\;\frac{R_k}{R_1+R_2+\dots+R_n}.
\]
\textbf{Divisor de corrente} (resistores em \emph{paralelo}): a corrente total se
reparte na proporção \emph{inversa} das resistências.
\[
I_k = I_{\text{total}}\;\frac{G_k}{G_1+\dots+G_n}
     = I_{\text{total}}\;\frac{1/R_k}{1/R_1+\dots+1/R_n}.
\]
\textbf{Caso de dois resistores} (o mais usado):
\[
U_1 = U\,\frac{R_1}{R_1+R_2}, \qquad
I_1 = I\,\frac{R_2}{R_1+R_2}\quad(\text{note o índice \emph{trocado}}).
\]
\end{conceito}

\begin{atencao}[O erro clássico]
No divisor de \textbf{corrente} de dois resistores, o ramo $R_1$ recebe a
corrente proporcional a $R_2$ (o \emph{outro} resistor). Trocar os índices faz a
maior corrente ir para o maior resistor --- o oposto do correto. Regra de ouro:
\emph{a maior corrente sempre passa pelo menor resistor}.
\end{atencao}

\legendaniveis

% =====================================================================
\subsection{Divisor de tensão}
% =====================================================================

\blocofixacao

\begin{questao}[1]{}
Em um divisor de tensão formado por dois resistores em série, a maior parcela de
tensão aparece:
\begin{alternativas}
\item sobre o menor resistor.
\item sobre o maior resistor.
\item igualmente sobre os dois, sempre.
\item sobre o resistor mais próximo do terminal positivo.
\item sobre nenhum: a tensão fica toda na fonte.
\end{alternativas}
\resposta{(b)}
\resolucao{Como a corrente é a mesma nos dois ($U = R\,I$), a tensão é
proporcional à resistência: o maior resistor fica com a maior tensão. A posição
em relação ao terminal positivo (alternativa d) é irrelevante.}
\end{questao}

\begin{questao}[1]{}
No circuito, três resistores em série de $\SI{3}{\ohm}$, $\SI{4}{\ohm}$ e
$\SI{5}{\ohm}$ estão ligados a uma fonte de $\SI{120}{\volt}$. Calcule a tensão
sobre cada resistor.

\circuitoq{%
\begin{circuitikz}[scale=0.95,transform shape,line width=0.8pt]
 \draw (0,0) to[battery1,l_={$E=\SI{120}{\volt}$}] (0,3)
       to[R,l={$R_1=\SI{3}{\ohm}$}] (3.2,3)
       to[R,l={$R_2=\SI{4}{\ohm}$}] (3.2,0)
       to[R,l={$R_3=\SI{5}{\ohm}$}] (0,0);
 \draw[->,Icol,line width=1pt] (0.7,3.4)--(1.7,3.4)
       node[midway,above,font=\footnotesize]{$I$};
\end{circuitikz}}

\resposta{$U_1=\SI{30}{\volt}$, $U_2=\SI{40}{\volt}$, $U_3=\SI{50}{\volt}$}
\resolucao{A corrente é $I = \dfrac{120}{3+4+5} = \dfrac{120}{12} = \SI{10}{\ampere}$.
Aplicando o divisor (ou simplesmente $U_k = R_k I$):
\[
U_1 = 3\cdot10 = \SI{30}{\volt},\quad
U_2 = 4\cdot10 = \SI{40}{\volt},\quad
U_3 = 5\cdot10 = \SI{50}{\volt}.
\]
\textbf{Verificação:} $30+40+50 = \SI{120}{\volt}$, igual à fonte.}
\end{questao}

\begin{questao}[1]{}
Uma fonte de $\SI{100}{\volt}$ alimenta dois resistores em série,
$R_1=\SI{60}{\ohm}$ e $R_2=\SI{40}{\ohm}$. A tensão de saída, tomada sobre $R_2$,
vale:
\begin{alternativas}[3]
\item \SI{20}{\volt}
\item \SI{40}{\volt}
\item \SI{50}{\volt}
\item \SI{60}{\volt}
\item \SI{100}{\volt}
\end{alternativas}
\resposta{(b)}
\resolucao{Aplicando diretamente a fórmula do divisor:
\[
U_2 = E\,\frac{R_2}{R_1+R_2} = 100\cdot\frac{40}{100} = \SI{40}{\volt}.
\]}
\end{questao}

\begin{questao}[1]{}
Três resistores iguais estão em série sob $\SI{60}{\volt}$. A tensão sobre cada um
vale:
\begin{alternativas}[3]
\item \SI{10}{\volt}
\item \SI{15}{\volt}
\item \SI{20}{\volt}
\item \SI{30}{\volt}
\item \SI{60}{\volt}
\end{alternativas}
\resposta{(c)}
\resolucao{Resistores iguais dividem a tensão igualmente:
$U = \dfrac{60}{3} = \SI{20}{\volt}$ em cada um. É o divisor de tensão mais
simples --- a base dos potenciômetros de precisão.}
\end{questao}

\blocoaplicacao

\begin{questao}[2]{}
Deseja-se obter $\SI{5}{\volt}$ a partir de uma fonte de $\SI{15}{\volt}$ usando
um divisor de dois resistores. Sabendo que $R_1 = \SI{2}{\kilo\ohm}$ (ligado ao
terminal positivo), determine $R_2$ (a saída é tomada sobre $R_2$).
\resposta{$R_2 = \SI{1}{\kilo\ohm}$}
\resolucao{Da relação do divisor,
\[
\frac{U_{\text{out}}}{E}=\frac{R_2}{R_1+R_2}
\ \Rightarrow\
\frac{5}{15}=\frac{R_2}{2000+R_2}
\ \Rightarrow\
2000+R_2 = 3R_2
\ \Rightarrow\
R_2 = \SI{1}{\kilo\ohm}.
\]
Repare que só importa a \emph{razão} $R_1{:}R_2 = 2{:}1$; qualquer par nessa
proporção daria \SI{5}{\volt} em vazio.}
\end{questao}

\begin{questao}[2]{}
No circuito da figura, calcule a tensão de saída $U_{\text{out}}$ (a) em vazio e
(b) quando uma carga $R_L = \SI{1}{\kilo\ohm}$ é ligada aos terminais de saída.
Comente o resultado.

\circuitoq{%
\begin{circuitikz}[scale=0.95,transform shape,line width=0.8pt]
 \draw (0,0) to[battery1,l_={$E=\SI{12}{\volt}$}] (0,3.4)
       to[R,l={$R_1=\SI{1}{\kilo\ohm}$}] (3.4,3.4) coordinate(a);
 \draw (a) to[R,l={$R_2=\SI{1}{\kilo\ohm}$}] (3.4,0) coordinate(b) -- (0,0);
 \draw (a) to[short,-o] (5,3.4) node[right]{$+$};
 \draw (b) to[short,-o] (5,0) node[right]{$-$};
 \draw[->] (5.3,3.0) -- (5.3,0.4) node[midway,right,font=\footnotesize]{$U_{\text{out}}$};
\end{circuitikz}}

\resposta{(a) \SI{6}{\volt}; (b) \SI{4}{\volt}}
\resolucao{(a) \textbf{Em vazio}, nenhuma corrente sai pelos terminais, e os dois
resistores iguais dividem a tensão ao meio:
$U_{\text{out}} = 12\cdot\dfrac{1}{2} = \SI{6}{\volt}$.\\
(b) \textbf{Com carga}, $R_L$ fica em paralelo com $R_2$:
\[
R_2 \parallel R_L = 1\text{k}\parallel 1\text{k} = \SI{0.5}{\kilo\ohm}
\ \Rightarrow\
U_{\text{out}} = 12\cdot\frac{0{,}5}{1+0{,}5} = \SI{4}{\volt}.
\]
\textbf{Conclusão:} a conexão da carga reduz a tensão de saída. Divisores resistivos só
mantêm a saída se $R_L \gg R_2$; por isso não servem como fonte de alimentação
para cargas variáveis --- para isso usa-se um regulador de tensão.}
\end{questao}

\begin{questao}[2]{}
Um potenciômetro de $\SI{10}{\kilo\ohm}$ é ligado como divisor de tensão a uma
fonte de $\SI{9}{\volt}$. Quando o cursor está a $\SI{30}{\percent}$ do curso (a
partir do terminal de referência), qual é a tensão de saída em vazio?
\resposta{\SI{2.7}{\volt}}
\resolucao{O cursor divide a resistência total na proporção do curso. A \SI{30}{\percent},
a saída corresponde a $\SI{30}{\percent}$ da tensão total:
\[
U_{\text{out}} = 9\cdot 0{,}30 = \SI{2.7}{\volt}.
\]
O valor da resistência total (\SI{10}{\kilo\ohm}) não entra na conta em vazio ---
apenas a \emph{fração} do curso importa. É exatamente assim que um controle de
volume ajusta o sinal.}
\end{questao}

\begin{questao}[2]{}
Para o circuito da figura (malha única com cinco resistores), calcule a corrente
e a tensão sobre o resistor de $\SI{8}{\ohm}$.

\circuitoq{%
\begin{circuitikz}[scale=0.9,transform shape,line width=0.8pt]
 \draw (0,0) to[battery1,l_={$E=\SI{100}{\volt}$},invert] (0,3)
       to[R,l={$\SI{10}{\ohm}$}] (3,3)
       to[R,l={$\SI{6}{\ohm}$}] (6,3)
       to[R,l={$\SI{8}{\ohm}$}] (6,0)
       to[R,l={$\SI{12}{\ohm}$}] (3,0)
       to[R,l={$\SI{6}{\ohm}$}] (0,0);
 \draw[->,Icol,line width=1pt] (0.7,3.4)--(1.7,3.4)
       node[midway,above,font=\footnotesize]{$I$};
\end{circuitikz}}

\resposta{$I = \SI{2.38}{\ampere}$ (aprox.); $U_{8} \approx \SI{19}{\volt}$}
\resolucao{Todos os resistores estão em série:
$\Req = 10+6+8+12+6 = \SI{42}{\ohm}$.
\[
I=\frac{100}{42}\approx\SI{2.38}{\ampere},
\qquad
U_8 = 8\cdot I = \frac{800}{42}\approx\SI{19.0}{\volt}.
\]
Pelo divisor: $U_8 = 100\cdot\dfrac{8}{42}\approx\SI{19.0}{\volt}$, o mesmo
resultado.}
\end{questao}

\blocoaprofundamento

\begin{questao}[3]{}
Um divisor de tensão deve fornecer $\SI{5}{\volt}$ a uma carga que consome
$\SI{10}{\milli\ampere}$, a partir de uma fonte de $\SI{15}{\volt}$. Como regra
de projeto, adota-se a \emph{corrente de sangria} (a que circula pelo divisor sem
passar pela carga) igual a $\SI{10}{\milli\ampere}$. Determine $R_1$ e $R_2$.

\circuitoq{%
\begin{circuitikz}[scale=0.95,transform shape,line width=0.8pt]
 \draw (0,0) to[battery1,l_={$E=\SI{15}{\volt}$}] (0,3.4)
       to[R,l={$R_1$},i={$I_1$}] (3.4,3.4) coordinate(a);
 \draw (a) to[R,l={$R_2$},i={$I_2$}] (3.4,0) coordinate(b) -- (0,0);
 \draw (a) to[short,-o] (5.2,3.4) coordinate(o1);
 \draw (b) to[short,-o] (5.2,0) coordinate(o2);
 \draw (o1) to[R,l={carga},i={$\SI{10}{\milli\ampere}$}] (o2);
\end{circuitikz}}

\resposta{$R_2 = \SI{500}{\ohm}$; $R_1 = \SI{500}{\ohm}$}
\resolucao{\textbf{Ramo de $R_2$ (sangria):} $R_2$ está em paralelo com a carga e
sob \SI{5}{\volt}. Com corrente de sangria $I_2 = \SI{10}{\milli\ampere}$:
\[
R_2 = \frac{U_2}{I_2}=\frac{5}{\pot{10}{-3}}=\SI{500}{\ohm}.
\]
\textbf{Ramo de $R_1$:} por ele passa a soma da corrente de sangria com a da carga,
$I_1 = 10+10 = \SI{20}{\milli\ampere}$, sob a tensão restante
$15-5 = \SI{10}{\volt}$:
\[
R_1=\frac{U_1}{I_1}=\frac{10}{\pot{20}{-3}}=\SI{500}{\ohm}.
\]
\textbf{Por que a sangria?} Se $R_2$ conduzisse muito pouco, a tensão de saída
variaria fortemente com a carga. Uma sangria comparável à corrente da carga
estabiliza a saída, ao custo de mais potência dissipada --- o compromisso típico
de projeto.}
\end{questao}

\begin{questao}[3]{}
Um divisor de tensão com $R_1 = \SI{4}{\kilo\ohm}$ e $R_2 = \SI{2}{\kilo\ohm}$ é
alimentado por $\SI{12}{\volt}$; a saída é tomada sobre $R_2$. Determine a faixa
de valores da carga $R_L$ para que a tensão de saída não caia mais que
$\SI{10}{\percent}$ em relação ao valor em vazio.
\resposta{$R_L \gtrsim \SI{12}{\kilo\ohm}$}
\resolucao{Em vazio, $U_0 = 12\cdot\dfrac{2}{6} = \SI{4}{\volt}$. Queremos
$U \geq 0{,}9\cdot U_0 = \SI{3.6}{\volt}$. Com carga, $R_2$ é substituído por
$R_2' = R_2\parallel R_L$:
\[
U = 12\cdot\frac{R_2'}{R_1+R_2'}\geq 3{,}6.
\]
Isolando, $\dfrac{R_2'}{4000+R_2'}\geq 0{,}3 \Rightarrow R_2'\geq \SI{1714}{\ohm}$
(aprox.). De $R_2' = \dfrac{2000\,R_L}{2000+R_L}\geq 1714$ resulta
$R_L \geq \SI{12}{\kilo\ohm}$ (aproximadamente). Ou seja, a carga precisa ser cerca
de \textbf{6 vezes maior} que $R_2$ para "não pesar" no divisor. Confirma a regra
prática: divisores resistivos só servem para cargas de alta impedância, como a
entrada de um amplificador.}
\end{questao}

\begin{questao}[3]{}
Um sistema de alimentação deve entregar $\SI{5}{\volt}$ a uma carga de
$\SI{100}{\ohm}$, a partir de uma fonte de $\SI{15}{\volt}$, usando um divisor
resistivo. Especifica-se que a tensão na carga não deve cair mais de
$\SI{2}{\percent}$ se a carga variar de $\SI{100}{\ohm}$ a $\SI{200}{\ohm}$.
\begin{itensq}
\item Explique por que um divisor puramente resistivo tem dificuldade em atender
      a essa especificação.
\item Determine a ordem de grandeza que $R_2$ (resistor em paralelo com a carga)
      deveria ter e discuta o custo dessa escolha.
\end{itensq}
\resposta{Ver análise; $R_2 \ll R_L$, com alto consumo}
\resolucao{(a) A tensão de saída de um divisor \emph{depende da carga}, pois
$R_L$ entra em paralelo com $R_2$. Quando $R_L$ varia, o paralelo
$R_2 \parallel R_L$ varia, e a saída varia junto. A regulação só é boa quando
$R_2 \parallel R_L \approx R_2$, isto é, quando $R_2 \ll R_L$.\\
(b) Quantificando: para que a variação de $R_L$ (de \SI{100}{} para
\SI{200}{\ohm}) altere o paralelo em menos de \SI{2}{\percent}, precisamos de
\[
R_2 \parallel 100 \approx R_2 \parallel 200
\quad\Longrightarrow\quad
R_2 \lesssim \SI{4}{\ohm} \ (\text{ordem de } R_L/25).
\]
Com $R_2 = \SI{4}{\ohm}$ e a razão do divisor exigindo
$R_1 = 2R_2 = \SI{8}{\ohm}$, a corrente do divisor seria
$I \approx \dfrac{15}{12} = \SI{1.25}{\ampere}$, dissipando cerca de
$\SI{19}{\watt}$ --- para entregar apenas $\dfrac{5^2}{100} = \SI{0.25}{\watt}$
à carga!\\
\textbf{Conclusão de projeto:} o rendimento seria de aproximadamente
\SI{1}{\percent}. Isso demonstra por que divisores resistivos servem apenas como
\emph{referência de tensão} para entradas de alta impedância, e por que fontes
reais usam \textbf{reguladores} (série ou chaveados), que mantêm a saída
independentemente da carga sem desperdiçar potência.}
\end{questao}

% BEGIN BANCO COMPLEMENTAR Divisor de tensão
% =====================================================================
%  Banco complementar --- Divisor de Tensao  (REVISADO)
%  Substitui a versao gerada por template (8 clones em N2, 9 em N3 e
%  8 questoes conceituais marcadas como N4).
%  O cap04 ja cobre: maior parcela (MC), tres resistores em serie,
%  100 V com 60/40, tres iguais sob 60 V, projeto de 5 V a partir de
%  15 V, saida em vazio e com carga R_L, potenciometro de 10 kohm como
%  divisor, malha unica com cinco resistores, divisor para carga de
%  10 mA, divisor 4k/2k carregado e sistema de 5 V para carga de
%  100 ohm. O EFEITO DE CARGA basico, portanto, ja esta no capitulo --
%  e tambem no banco de circuitos mistos. Aqui ele so reaparece na
%  forma da CONDICAO DE PROJETO (N4) e da nao linearidade do
%  potenciometro carregado, que sao angulos novos.
%  Distribuicao mantida: 4 N1 + 8 N2 + 9 N3 + 8 N4 = 29
% =====================================================================

\subsubsection*{Banco complementar --- Divisor de tensão}

\begin{atencao}[Organização do banco]
A relação $U_o=U_{in}\dfrac{R_2}{R_1+R_2}$ pressupõe \textbf{corrente comum} ---
ou seja, saída em vazio. Todo o resto do tópico decorre de levar a sério essa
hipótese: o divisor tem resistência interna $R_1\parallel R_2$, e qualquer carga
comparável a ela derruba a saída. O N3 explora tomadas múltiplas, sensores e
fontes não ideais; o N4 trata de projeto sob restrição de potência, tolerância
de razão, sonda compensada e a comparação com reguladores.
\end{atencao}

\blocofixacao

\begin{questao}[1]{}
Três resistores de $\SI{1}{\kilo\ohm}$, $\SI{2}{\kilo\ohm}$ e
$\SI{3}{\kilo\ohm}$ estão em série sob $\SI{12}{\volt}$. Determine a tensão
sobre o resistor \emph{intermediário}.
\resposta{$U=\SI{4}{\volt}$}
\resolucao{A regra do divisor vale para qualquer número de elementos em série:
cada um fica com a fração de sua própria resistência sobre o total.
\[
U_2=12\cdot\frac{2}{1+2+3}=12\cdot\frac13=\SI{4}{\volt}.
\]
\textbf{Conferência:} as três tensões são $2$, $4$ e $\SI{6}{\volt}$, somando
$\SI{12}{\volt}$ \checkmark, e guardam a proporção $1{:}2{:}3$ dos
resistores.}
\end{questao}

\begin{questao}[1]{}
Um divisor deve entregar $25\%$ da tensão de entrada. Determine a razão
$R_1/R_2$.
\resposta{$R_1/R_2=3$}
\resolucao{
\[
\frac{R_2}{R_1+R_2}=0{,}25
\;\Longrightarrow\;
R_1+R_2=4R_2
\;\Longrightarrow\;
R_1=3R_2.
\]
\textbf{Observação importante:} apenas a \emph{razão} é fixada --- $3\,$k$/1\,$k
e $30\,$k$/10\,$k produzem a mesma saída em vazio. O que os distingue é a
corrente de repouso e a resistência interna do divisor, e é isso que decide o
projeto quando há carga.}
\end{questao}

\begin{questao}[1]{}
Em um divisor em vazio, ambos os resistores são \textbf{dobrados}. A tensão de
saída:
\begin{alternativas}[2]
\item dobra
\item cai à metade
\item não se altera
\item depende dos valores originais
\end{alternativas}
\resposta{(c)}
\resolucao{A razão $\dfrac{2R_2}{2R_1+2R_2}=\dfrac{R_2}{R_1+R_2}$ é
inalterada --- a saída em vazio depende só da proporção.\\
\textbf{O que muda:} a corrente de repouso cai à metade (menos consumo) e a
resistência interna $R_1\parallel R_2$ \emph{dobra} (mais sensível à carga).
Escalar um divisor troca consumo por robustez, e essa é sempre a decisão
central do projeto.}
\end{questao}

\begin{questao}[1]{}
Analise os casos-limite de um divisor: (a) $R_2\to0$; (b) $R_2\to\infty$.
\resposta{(a) $U_o\to0$; (b) $U_o\to U_{in}$}
\resolucao{(a) Com $R_2=0$, a saída está em \textbf{curto} com o terra:
$U_o=0$. A corrente vale $U_{in}/R_1$ --- toda a potência vai para $R_1$.\\
(b) Com $R_2\to\infty$ (circuito aberto), não circula corrente, não há queda em
$R_1$ e $U_o\to U_{in}$.\\
\textbf{Conclusão:} a saída de um divisor está sempre entre $0$ e $U_{in}$ --- ele
só \emph{reduz} tensão, nunca a eleva. Obter tensão maior que a entrada exige
elemento armazenador (indutor ou capacitor) e comutação.}
\end{questao}

\blocoaplicacao

\begin{questao}[2]{}
Um divisor de $\SI{12}{\volt}$ usa três resistores em série de
$\SI{1}{\kilo\ohm}$, $\SI{1}{\kilo\ohm}$ e $\SI{2}{\kilo\ohm}$, com duas
tomadas intermediárias.
\begin{itensq}
\item Determine as duas tensões de saída em relação ao terra.
\item Determine a corrente de repouso.
\end{itensq}
\resposta{(a) $\SI{9}{\volt}$ e $\SI{6}{\volt}$; (b) $\SI{3}{\milli\ampere}$}
\resolucao{Total: $R=\SI{4}{\kilo\ohm}$ e
$I=12/4000=\SI{3}{\milli\ampere}$.\\
A tomada mais alta tem, abaixo de si, $1+2=\SI{3}{\kilo\ohm}$:
\[
U_A=12\cdot\frac{3}{4}=\SI{9}{\volt}.
\]
A tomada inferior tem $\SI{2}{\kilo\ohm}$ abaixo:
\[
U_B=12\cdot\frac{2}{4}=\SI{6}{\volt}.
\]
\textbf{Cuidado:} cada saída conta apenas a resistência \emph{entre ela e o
terra}, não a do trecho acima. E as duas tomadas se \emph{influenciam}: carregar
uma altera a outra, pois compartilham a mesma cadeia.}
\end{questao}

\begin{questao}[2]{}
Uma fonte de $\SI{30}{\volt}$ alimenta dois resistores iguais em série, e o
ponto médio é adotado como referência (terra). Determine as tensões dos dois
extremos.
\resposta{$\SI{+15}{\volt}$ e $\SI{-15}{\volt}$}
\resolucao{Cada resistor cai $\SI{15}{\volt}$. Tomando o ponto médio como
$\SI{0}{\volt}$, o terminal superior fica em $\SI{+15}{\volt}$ e o inferior em
$\SI{-15}{\volt}$.\\
\textbf{Aplicação:} é assim que se obtém alimentação \emph{simétrica} para
amplificadores operacionais a partir de uma fonte simples. \\
\textbf{Limitação séria:} o "terra" assim criado é um ponto de divisor, com
resistência interna $R/2$. Qualquer corrente drenada de forma assimétrica pelos
dois ramos desloca a referência. Na prática, o divisor precisa ser reforçado por
um seguidor de tensão ou por capacitores de grande valor.}
\end{questao}

\begin{questao}[2]{}
Um divisor de $\SI{24}{\volt}$ tem $R_1=\SI{8}{\kilo\ohm}$ e
$R_2=\SI{4}{\kilo\ohm}$. Determine a saída, a corrente de repouso e a potência
total dissipada.
\resposta{$\SI{8}{\volt}$; $\SI{2}{\milli\ampere}$; $\SI{48}{\milli\watt}$}
\resolucao{
\[
U_o=24\cdot\frac{4}{12}=\SI{8}{\volt};
\qquad
I=\frac{24}{12000}=\SI{2}{\milli\ampere};
\]
\[
P=U_{in}I=24(\pot{2}{-3})=\SI{48}{\milli\watt}.
\]
\textbf{Distribuição:} $P_1=8000(\pot{2}{-3})^{2}=\SI{32}{\milli\watt}$ e
$P_2=\SI{16}{\milli\watt}$ --- proporcionais às resistências, como em qualquer
série.\\
Note que essa potência é consumida \emph{permanentemente}, mesmo sem carga
alguma conectada. Em equipamentos alimentados por bateria, divisores de repouso
são uma fonte silenciosa de descarga.}
\end{questao}

\begin{questao}[2]{}
Um divisor tem $R_1=\SI{2}{\kilo\ohm}$ fixo e $R_2$ ajustável de
$\SI{0}{\ohm}$ a $\SI{8}{\kilo\ohm}$, sob $\SI{10}{\volt}$. Determine a faixa
da tensão de saída.
\resposta{de $\SI{0}{\volt}$ a $\SI{8}{\volt}$}
\resolucao{
\[
R_2=0 \Rightarrow U_o=0;
\qquad
R_2=\SI{8}{\kilo\ohm} \Rightarrow U_o=10\cdot\frac{8}{10}=\SI{8}{\volt}.
\]
\textbf{Não linearidade do ajuste:} no meio do curso
($R_2=\SI{4}{\kilo\ohm}$), a saída é $10\cdot4/6=\SI{6.67}{\volt}$ --- e não
$\SI{4}{\volt}$, como uma escala linear sugeriria. A relação
$U_o(R_2)$ é uma hipérbole, não uma reta.\\
Para obter ajuste \emph{linear}, é preciso usar um potenciômetro (em que a soma
$R_1+R_2$ permanece constante), e não um reostato em série com resistor fixo.}
\end{questao}

\begin{questao}[2]{}
Um divisor deve entregar $\SI{4.5}{\volt}$ com $R_2=\SI{3}{\kilo\ohm}$.
Determine $U_{in}$ para $R_1=\SI{5}{\kilo\ohm}$.
\resposta{$U_{in}=\SI{12}{\volt}$}
\resolucao{
\[
4{,}5=U_{in}\cdot\frac{3}{5+3}
\;\Longrightarrow\;
U_{in}=4{,}5\cdot\frac{8}{3}=\SI{12}{\volt}.
\]
\textbf{Verificação:} $12\cdot3/8=\SI{4.5}{\volt}$ \checkmark, e a corrente é
$12/8000=\SI{1.5}{\milli\ampere}$.\\
A relação do divisor pode ser resolvida em qualquer das quatro variáveis ---
$U_{in}$, $U_o$, $R_1$ ou $R_2$ --- desde que as outras três sejam conhecidas.}
\end{questao}

\begin{questao}[2]{}
Um termistor NTC vale $\SI{10}{\kilo\ohm}$ a $\SI{25}{\celsius}$ e
$\SI{4}{\kilo\ohm}$ a $\SI{50}{\celsius}$. Ele ocupa a posição de $R_2$ em um
divisor com $R_1=\SI{10}{\kilo\ohm}$ sob $\SI{5}{\volt}$. Determine a saída nas
duas temperaturas.
\resposta{$\SI{2.5}{\volt}$ e $\SI{1.43}{\volt}$}
\resolucao{
\[
\SI{25}{\celsius}:\ U_o=5\cdot\frac{10}{20}=\SI{2.5}{\volt};
\qquad
\SI{50}{\celsius}:\ U_o=5\cdot\frac{4}{14}=\SI{1.43}{\volt}.
\]
Variação de $\SI{1.07}{\volt}$ em $\SI{25}{\celsius}$, ou cerca de
$\SI{43}{\milli\volt\per\celsius}$ na média --- sinal amplo e fácil de digitalizar.\\
\textbf{Ressalva:} a relação \emph{não} é linear, pois nem o termistor é linear
em $T$, nem o divisor é linear em $R_2$. Para converter tensão em temperatura é
preciso tabela, polinômio ou a equação de Steinhart--Hart. A escolha de
$R_1=\SI{10}{\kilo\ohm}$ não é casual --- ela maximiza a sensibilidade em torno
de $\SI{25}{\celsius}$, como se demonstra no bloco de desafio.}
\end{questao}

\begin{questao}[2]{}
Um LDR varia de $\SI{1}{\kilo\ohm}$ (claro) a $\SI{100}{\kilo\ohm}$ (escuro).
Ele é colocado como $R_1$ de um divisor com $R_2=\SI{10}{\kilo\ohm}$ sob
$\SI{5}{\volt}$. Determine a faixa de saída e o sentido da variação.
\resposta{de $\SI{4.55}{\volt}$ (claro) a $\SI{0.45}{\volt}$ (escuro)}
\resolucao{
\[
\text{claro}:\ U_o=5\cdot\frac{10}{11}=\SI{4.55}{\volt};
\qquad
\text{escuro}:\ U_o=5\cdot\frac{10}{110}=\SI{0.45}{\volt}.
\]
Com o LDR na posição \emph{superior}, mais luz significa \emph{maior} tensão de
saída.\\
\textbf{Inversão do comportamento:} trocando as posições (LDR embaixo), a saída
passaria de $\SI{0.45}{\volt}$ no claro a $\SI{4.55}{\volt}$ no escuro. A
escolha da posição define a polaridade do sensor --- decisão de projeto que
custa nada e evita inverter o sinal depois em software.}
\end{questao}

\begin{questao}[2]{}
Um divisor projetado com $R_1=\SI{8.7}{\kilo\ohm}$ e $R_2=\SI{3.3}{\kilo\ohm}$
para $12\to\SI{3.3}{\volt}$ deve usar valores E24. Adotando
$R_1=\SI{9.1}{\kilo\ohm}$, determine a saída real e o erro.
\resposta{$U_o=\SI{3.19}{\volt}$; erro de $-3{,}2\%$}
\resolucao{
\[
U_o=12\cdot\frac{3{,}3}{9{,}1+3{,}3}=12\cdot\frac{3{,}3}{12{,}4}=\SI{3.19}{\volt};
\qquad
\text{erro}=\frac{3{,}19-3{,}30}{3{,}30}=-3{,}2\%.
\]
\textbf{Alternativa melhor:} o par $\SI{10}{\kilo\ohm}/\SI{3.9}{\kilo\ohm}$ dá
$U_o=\SI{3.37}{\volt}$, erro de apenas $+2{,}0\%$. Vale sempre testar
\emph{vários} pares E24 antes de escolher --- a razão importa mais que os
valores individuais.\\
\textbf{Se $3\%$ for demais:} use a série E96 (tolerância $1\%$), que contém
$\SI{8.66}{\kilo\ohm}$ e resolve o problema diretamente.}
\end{questao}

\blocoaprofundamento

\begin{questao}[3]{}
Um divisor $R_1=\SI{6}{\kilo\ohm}$, $R_2=\SI{3}{\kilo\ohm}$ é alimentado por
uma fonte \textbf{real} de $\SI{12}{\volt}$ com $r=\SI{1}{\kilo\ohm}$.
\begin{itensq}
\item Determine a saída, considerando $r$.
\item Compare com a previsão que ignora $r$.
\item Determine a resistência de Thévenin vista na saída.
\end{itensq}
\resposta{(a) $\SI{3.6}{\volt}$; (b) previsão de $\SI{4}{\volt}$, erro de
$10\%$; (c) $R_{Th}=\SI{2.1}{\kilo\ohm}$}
\resolucao{(a) A resistência interna soma-se a $R_1$:
\[
U_o=12\cdot\frac{3}{1+6+3}=12\cdot\frac{3}{10}=\SI{3.6}{\volt}.
\]
(b) Ignorando $r$: $12\cdot3/9=\SI{4}{\volt}$. O erro é de $10\%$ --- nada
desprezível.\\
(c) Zerando a f.e.m. (mas mantendo $r$), a saída enxerga $R_2$ em paralelo com
$(r+R_1)$:
\[
R_{Th}=\frac{3\cdot7}{10}=\SI{2.1}{\kilo\ohm}.
\]
\textbf{Leitura:} a resistência interna da fonte se comporta como parte de
$R_1$ --- ela piora a razão \emph{e} aumenta a impedância de saída. Em divisores
de baixa impedância (resistores de dezenas de ohms), $r$ pode dominar
completamente o projeto.}
\end{questao}

\begin{questao}[3]{}
Dois divisores idênticos ($R_1=\SI{3}{\kilo\ohm}$, $R_2=\SI{1}{\kilo\ohm}$)
são alimentados pela mesma fonte de $\SI{8}{\volt}$, mas em um deles $R_2$ é
trocado por $\SI{1.1}{\kilo\ohm}$.
\begin{itensq}
\item Determine as duas saídas e a diferença entre elas.
\item Determine a sensibilidade da diferença a variações de $R_2$.
\end{itensq}
\resposta{(a) $\SI{2}{\volt}$ e $\SI{2.146}{\volt}$; diferença de
$\SI{146}{\milli\volt}$}
\resolucao{(a)
\[
U_A=8\cdot\frac{1}{4}=\SI{2}{\volt};
\qquad
U_B=8\cdot\frac{1{,}1}{4{,}1}=\SI{2.146}{\volt};
\qquad
\Delta U=\SI{146}{\milli\volt}.
\]
(b) Derivando $U=U_{in}R_2/(R_1+R_2)$:
\[
\frac{\mathrm{d}U}{\mathrm{d}R_2}
=\frac{U_{in}R_1}{(R_1+R_2)^{2}}
=\frac{8(3000)}{(4000)^{2}}=\SI{1.5}{\milli\volt\per\ohm}.
\]
Para $\Delta R_2=\SI{100}{\ohm}$, prevê-se $\SI{150}{\milli\volt}$ --- muito
próximo dos $\SI{146}{\milli\volt}$ exatos, pois a variação é pequena.\\
\textbf{Esta é a estrutura da ponte:} dois divisores comparados, em que o
\emph{desequilíbrio} carrega a informação. A grande vantagem é que variações
comuns aos dois ramos --- deriva da fonte, temperatura ambiente --- se cancelam
na diferença.}
\end{questao}

\begin{questao}[3]{}
Um divisor $R_1=R_2=\SI{100}{\kilo\ohm}$ sob $\SI{10}{\volt}$ é medido por um
voltímetro cuja resistência de entrada é $\SI{1}{\mega\ohm}$.
\begin{itensq}
\item Determine a leitura do instrumento.
\item Determine o erro percentual.
\item Determine a resistência mínima do voltímetro para erro inferior a
      $1\%$.
\end{itensq}
\resposta{(a) $\SI{4.76}{\volt}$; (b) $-4{,}8\%$; (c)
$\ge\SI{5}{\mega\ohm}$}
\resolucao{(a) O voltímetro entra em paralelo com $R_2$:
\[
R_2'=\frac{100\cdot1000}{1100}=\SI{90.9}{\kilo\ohm};
\qquad
U=10\cdot\frac{90{,}9}{100+90{,}9}=\SI{4.76}{\volt}.
\]
(b) Erro: $(4{,}76-5{,}00)/5{,}00=-4{,}8\%$.\\
(c) A resistência interna do divisor é
$R_{Th}=100\parallel100=\SI{50}{\kilo\ohm}$. O erro relativo é
aproximadamente $R_{Th}/R_V$, então
\[
\frac{50\,\mathrm{k}}{R_V}\le0{,}01
\;\Longrightarrow\;
R_V\ge\SI{5}{\mega\ohm}.
\]
\textbf{Lição de instrumentação:} o instrumento \emph{perturba} o circuito que
mede. Multímetros digitais comuns têm $\SI{10}{\mega\ohm}$ --- suficientes aqui,
mas insuficientes para divisores de megaohms. E note que a especificação
depende de $R_{Th}$ do circuito, não da tensão medida.}
\end{questao}

\begin{questao}[3]{}
Um divisor $R_1=\SI{4}{\kilo\ohm}$, $R_2=\SI{1}{\kilo\ohm}$ sob
$\SI{10}{\volt}$ fornece referência a um conversor A/D de $10$ bits com fundo
de escala de $\SI{2.5}{\volt}$.
\begin{itensq}
\item Determine a saída e a margem em relação ao fundo de escala.
\item Determine o degrau de quantização.
\item Compare o erro de quantização com o erro de resistores de $1\%$.
\end{itensq}
\resposta{(a) $\SI{2}{\volt}$, $80\%$ do fundo; (b) $\SI{2.44}{\milli\volt}$;
(c) resistores dominam ($\approx\SI{16}{\milli\volt}$)}
\resolucao{(a) $U_o=10\cdot1/5=\SI{2}{\volt}$, ou $80\%$ do fundo de escala ---
boa utilização da faixa, sem risco de saturação.\\
(b) Com $10$ bits, $2^{10}=1024$ degraus:
\[
\Delta=\frac{2{,}5}{1024}=\SI{2.44}{\milli\volt}.
\]
(c) Com resistores de $1\%$, o pior erro de razão vale
$\dfrac{R_1}{R_1+R_2}(t_1+t_2)=0{,}8(2\%)=1{,}6\%$, ou
$0{,}016\times2=\SI{32}{\milli\volt}$ --- \textbf{treze vezes} o degrau de
quantização.\\
\textbf{Conclusão de projeto:} não adianta usar um conversor de $12$ ou
$16$ bits com esse divisor. O gargalo é o divisor, e a solução é resistores de
$0{,}1\%$, calibração por software, ou uma referência de tensão dedicada.
Aumentar a resolução do conversor sem corrigir o divisor é gastar dinheiro em
dígitos sem significado.}
\end{questao}

\begin{questao}[3]{}
Um divisor de $\SI{48}{\volt}$ para $\SI{12}{\volt}$ deve alimentar uma carga
de $\SI{100}{\ohm}$.
\begin{itensq}
\item Determine $R_1$ e $R_2$ ignorando a carga, com corrente de repouso de
      $\SI{120}{\milli\ampere}$.
\item Determine a saída real com a carga conectada.
\item Avalie a viabilidade.
\end{itensq}
\resposta{(a) $R_1=\SI{300}{\ohm}$, $R_2=\SI{100}{\ohm}$;
(b) $\SI{6.86}{\volt}$; (c) inviável}
\resolucao{(a) $R_2=12/0{,}120=\SI{100}{\ohm}$;
$R_1=(48-12)/0{,}120=\SI{300}{\ohm}$.\\
(b) A carga de $\SI{100}{\ohm}$ fica em paralelo com $R_2$:
\[
R_2'=\frac{100\cdot100}{200}=\SI{50}{\ohm};
\qquad
U_o=48\cdot\frac{50}{350}=\SI{6.86}{\volt}.
\]
Queda de $43\%$ em relação ao projetado.\\
(c) \textbf{Inviável.} A carga consome $\SI{120}{\milli\ampere}$ a
$\SI{12}{\volt}$ --- exatamente a corrente de repouso, e não uma fração dela. A
regra de bolso exige repouso de $10$ a $20$ vezes a corrente da carga, o que
levaria $R_2$ a $\SI{10}{\ohm}$ e a dissipação total a
$48\times1{,}2=\SI{57.6}{\watt}$ para entregar $\SI{1.44}{\watt}$ --- rendimento
de $2{,}5\%$.\\
\textbf{Diagnóstico geral:} divisores resistivos servem para gerar
\emph{referências} de alta impedância, não para \emph{alimentar} cargas. Aqui o
caso pede regulador.}
\end{questao}

\begin{questao}[3]{}
Um divisor com $R_1=\SI{10}{\kilo\ohm}$ e $R_2=\SI{10}{\kilo\ohm}$ opera sob
$\SI{200}{\volt}$.
\begin{itensq}
\item Determine a potência de cada resistor e a tensão sobre cada um.
\item Especifique os resistores considerando potência \emph{e} tensão máxima de
      trabalho.
\end{itensq}
\resposta{$\SI{1}{\watt}$ e $\SI{100}{\volt}$ em cada; exige resistor de
potência com tensão de trabalho adequada}
\resolucao{(a) $I=200/20000=\SI{10}{\milli\ampere}$;
\[
U_1=U_2=\SI{100}{\volt};
\qquad
P_1=P_2=10^{4}(10^{-2})^{2}=\SI{1}{\watt}.
\]
(b) \textbf{Dois critérios independentes:}
\begin{itemize}
\item \textbf{Potência:} com fator $2$ de folga, resistores de
$\SI{2}{\watt}$.
\item \textbf{Tensão máxima de trabalho:} um resistor de filme comum de
$\frac14\,\si{\watt}$ suporta cerca de $\SI{250}{\volt}$; os de
$\SI{2}{\watt}$ chegam a $\SI{500}{\volt}$. Os $\SI{100}{\volt}$ exigidos são
atendidos com folga.
\end{itemize}
\textbf{Por que o segundo critério existe:} a tensão máxima é limitada pela
ruptura do revestimento e pelo arco entre as espiras do filme --- fenômeno que
independe da potência. Em divisores de \emph{alta} tensão (quilovolts), é comum
que a tensão, e não a potência, force o uso de vários resistores em série,
mesmo com dissipação irrisória.}
\end{questao}

\begin{questao}[3]{}
Um divisor deve fornecer $\SI{5}{\volt}$ a partir de $\SI{15}{\volt}$ para uma
entrada de microcontrolador com corrente de entrada de
$\SI{1}{\micro\ampere}$.
\begin{itensq}
\item Determine $R_1$ e $R_2$ para corrente de repouso de
      $\SI{10}{\micro\ampere}$.
\item Verifique o erro introduzido pela corrente de entrada.
\item Comente o compromisso com o consumo.
\end{itensq}
\resposta{(a) $R_1=\SI{1}{\mega\ohm}$, $R_2=\SI{500}{\kilo\ohm}$;
(b) erro de $\approx3\%$}
\resolucao{(a) $R_2=5/\pot{10}{-6}=\SI{500}{\kilo\ohm}$;
$R_1=10/\pot{10}{-6}=\SI{1}{\mega\ohm}$.\\
(b) A corrente de entrada de $\SI{1}{\micro\ampere}$ é $10\%$ da corrente de
repouso --- longe de desprezível. Modelando a entrada como resistência
equivalente $R_{in}=5/\pot{1}{-6}=\SI{5}{\mega\ohm}$:
\[
R_2'=\frac{500\cdot5000}{5500}=\SI{454.5}{\kilo\ohm};
\qquad
U_o=15\cdot\frac{454{,}5}{1454{,}5}=\SI{4.69}{\volt},
\]
erro de $-6{,}2\%$.\\
(c) \textbf{Compromisso.} Reduzir o erro exige aumentar a corrente de repouso ---
com $\SI{100}{\micro\ampere}$ ($R_1=\SI{100}{\kilo\ohm}$,
$R_2=\SI{50}{\kilo\ohm}$), o erro cai para cerca de $0{,}7\%$, mas o consumo em
repouso sobe dez vezes.\\
Em um dispositivo alimentado por bateria de $\SI{1000}{\milli\ampere\hour}$,
$\SI{100}{\micro\ampere}$ contínuos consomem a bateria em pouco mais de um ano
\emph{só com o divisor}. É o tipo de detalhe que decide a autonomia de um
produto.}
\end{questao}

\begin{questao}[3]{}
Determine, em um divisor, a expressão da resistência de Thévenin vista na saída
e avalie-a para $R_1=\SI{6}{\kilo\ohm}$ e $R_2=\SI{3}{\kilo\ohm}$.
\begin{itensq}
\item Deduza $R_{Th}$.
\item Mostre que $R_{Th}<\min(R_1,R_2)$.
\item Explique seu papel no carregamento.
\end{itensq}
\resposta{$R_{Th}=R_1\parallel R_2=\SI{2}{\kilo\ohm}$}
\resolucao{(a) Zerando a fonte de entrada (curto), a saída enxerga $R_1$ e
$R_2$ em paralelo:
\[
R_{Th}=\frac{R_1R_2}{R_1+R_2}=\frac{6\cdot3}{9}=\SI{2}{\kilo\ohm}.
\]
(b) O paralelo de dois resistores positivos é sempre menor que o menor deles ---
resultado geral já estabelecido para associações. Aqui,
$\SI{2}{\kilo\ohm}<\SI{3}{\kilo\ohm}$ \checkmark.\\
(c) \textbf{Papel no carregamento.} Conectada uma carga $R_L$, o divisor se
comporta como fonte de Thévenin:
\[
U_L=U_{Th}\cdot\frac{R_L}{R_L+R_{Th}}.
\]
O erro relativo é aproximadamente $R_{Th}/R_L$ para $R_L\gg R_{Th}$. Logo:
\begin{itemize}
\item $R_L=10R_{Th}$ $\Rightarrow$ erro $\approx9\%$;
\item $R_L=100R_{Th}$ $\Rightarrow$ erro $\approx1\%$.
\end{itemize}
\textbf{Consequência de projeto:} o único parâmetro que importa para o
carregamento é $R_{Th}$ --- não $R_1$, não $R_2$, não a razão. Reduzir os dois
resistores proporcionalmente reduz $R_{Th}$ e melhora a robustez, ao custo de
consumo.}
\end{questao}

\begin{questao}[3]{}
Um divisor alimenta uma carga que \textbf{varia} entre $\SI{10}{\kilo\ohm}$ e
$\SI{100}{\kilo\ohm}$. O divisor tem $R_{Th}=\SI{2}{\kilo\ohm}$ e
$U_{Th}=\SI{5}{\volt}$.
\begin{itensq}
\item Determine a faixa de tensão de saída.
\item Determine a regulação percentual.
\item Indique como melhorá-la.
\end{itensq}
\resposta{(a) de $\SI{4.17}{\volt}$ a $\SI{4.90}{\volt}$;
(b) $\approx15\%$}
\resolucao{(a)
\[
R_L=\SI{10}{\kilo\ohm}:\ U=5\cdot\frac{10}{12}=\SI{4.17}{\volt};
\qquad
R_L=\SI{100}{\kilo\ohm}:\ U=5\cdot\frac{100}{102}=\SI{4.90}{\volt}.
\]
(b) Regulação:
\[
\frac{4{,}90-4{,}17}{4{,}90}=15{,}0\%.
\]
(c) \textbf{Como melhorar.}
\begin{itemize}
\item \textbf{Reduzir $R_{Th}$.} Com $R_{Th}=\SI{200}{\ohm}$, a faixa passa a
$4{,}90$--$\SI{4.99}{\volt}$ e a regulação cai para $1{,}8\%$ --- ao custo de
dez vezes mais consumo de repouso.
\item \textbf{Inserir um seguidor de tensão} (amp-op em configuração de ganho
unitário). A impedância de saída cai para frações de ohm e a regulação
praticamente desaparece, sem aumentar o consumo do divisor.
\item \textbf{Usar regulador}, se a carga também exigir corrente apreciável.
\end{itemize}
\textbf{Observação:} a segunda opção é a mais elegante --- ela desacopla
\emph{razão} de \emph{impedância}, que no divisor puro estão inevitavelmente
amarradas.}
\end{questao}

\blocodesafio

\begin{questao}[4]{}
\textbf{(Condição de carregamento.)} Deduza a condição para que uma carga
$R_L$ altere a saída de um divisor em menos de um erro relativo $\varepsilon$.
\begin{itensq}
\item Obtenha a expressão exata do erro.
\item Deduza a condição aproximada para $\varepsilon$ pequeno.
\item Construa a tabela de $R_L/R_{Th}$ para
      $\varepsilon=10\%$, $1\%$ e $0{,}1\%$.
\end{itensq}
\resposta{$\varepsilon=\dfrac{R_{Th}}{R_L+R_{Th}}$; aproximadamente
$R_L\ge R_{Th}/\varepsilon$}
\resolucao{(a) Com carga,
$U_L=U_{Th}\dfrac{R_L}{R_L+R_{Th}}$, logo
\[
\varepsilon=\frac{U_{Th}-U_L}{U_{Th}}=1-\frac{R_L}{R_L+R_{Th}}
=\frac{R_{Th}}{R_L+R_{Th}}.
\]
(b) Para $R_L\gg R_{Th}$, o denominador tende a $R_L$:
\[
\varepsilon\approx\frac{R_{Th}}{R_L}
\;\Longrightarrow\;
R_L\ge\frac{R_{Th}}{\varepsilon}.
\]
(c)
\begin{center}
\begin{tabular}{ccc}
\hline
$\varepsilon$ & $R_L/R_{Th}$ (exato) & $R_L/R_{Th}$ (aprox.)\\
\hline
$10\%$ & $9$ & $10$\\
$1\%$ & $99$ & $100$\\
$0{,}1\%$ & $999$ & $1000$\\
\hline
\end{tabular}
\end{center}
A aproximação erra por exatamente $1$ unidade e é conservadora --- pode ser
usada sem receio.\\
\textbf{Regra de bolso resultante:} para $1\%$ de erro,
$R_L\ge100\,R_{Th}$. Note que $R_{Th}=R_1\parallel R_2$ é sempre \emph{menor}
que qualquer um dos resistores, o que torna a condição menos severa do que
comparar $R_L$ com $R_2$ diretamente --- erro comum que leva a
superdimensionar o divisor.}
\end{questao}

\begin{questao}[4]{}
\textbf{(Projeto sob restrição de potência.)} Projete um divisor que reduza
$\SI{100}{\volt}$ a $\SI{5}{\volt}$, dissipando no máximo
$\SI{0.5}{\watt}$.
\begin{itensq}
\item Determine $R_1$ e $R_2$.
\item Determine a potência e a tensão de cada resistor.
\item Determine a maior carga que pode ser atendida com erro de até $1\%$.
\end{itensq}
\resposta{(a) $R_1=\SI{19}{\kilo\ohm}$, $R_2=\SI{1}{\kilo\ohm}$;
(b) $0{,}475$ e $\SI{0.025}{\watt}$; (c) $R_L\ge\SI{95}{\kilo\ohm}$}
\resolucao{(a) A restrição de potência fixa a \emph{soma}:
\[
P=\frac{U_{in}^{2}}{R_1+R_2}\le0{,}5
\;\Longrightarrow\;
R_1+R_2\ge\frac{100^{2}}{0{,}5}=\SI{20}{\kilo\ohm}.
\]
A razão fixa a \emph{proporção}: $R_2/(R_1+R_2)=0{,}05$. Adotando o mínimo,
\[
R_2=\SI{1}{\kilo\ohm};\qquad R_1=\SI{19}{\kilo\ohm}.
\]
(b) $I=100/20000=\SI{5}{\milli\ampere}$;
\[
P_1=19000(\pot{5}{-3})^{2}=\SI{0.475}{\watt};
\qquad
P_2=\SI{0.025}{\watt};
\]
com $U_1=\SI{95}{\volt}$ e $U_2=\SI{5}{\volt}$.\\
\textbf{Especificação:} $R_1$ concentra $95\%$ da dissipação --- use
$\SI{1}{\watt}$ nele e $\frac18\,\si{\watt}$ em $R_2$. Alternativamente,
dividir $R_1$ em dois de $\SI{9.5}{\kilo\ohm}$ reparte o calor e a tensão.\\
(c) $R_{Th}=19\parallel1=\SI{950}{\ohm}$; pela condição anterior,
\[
R_L\ge100\,R_{Th}=\SI{95}{\kilo\ohm}.
\]
\textbf{Tensão de projeto:} o divisor entrega $\SI{5}{\volt}$ a uma carga que
consome no máximo $\SI{53}{\micro\ampere}$ --- serve para um conversor A/D, não
para alimentar nada.}
\end{questao}

\begin{questao}[4]{}
\textbf{(Potenciômetro carregado.)} Um potenciômetro de $\SI{10}{\kilo\ohm}$
é usado como divisor sob $\SI{10}{\volt}$, com carga
$R_L=\SI{10}{\kilo\ohm}$ na saída. Seja $x$ a fração do curso.
\begin{itensq}
\item Deduza $U_o(x)$ com a carga.
\item Avalie em $x=0{,}25$, $0{,}5$ e $0{,}75$ e compare com o caso ideal.
\item Determine onde o desvio é máximo e como reduzi-lo.
\end{itensq}
\resposta{Em $x=0{,}5$: $\SI{4.00}{\volt}$ contra $\SI{5}{\volt}$ ideal ---
desvio de $20\%$}
\resolucao{(a) A porção inferior vale $xR_p$ e fica em paralelo com $R_L$; a
superior vale $(1-x)R_p$:
\[
U_o=U_{in}\,\frac{xR_p\parallel R_L}{(1-x)R_p+\left(xR_p\parallel R_L\right)}.
\]
Com $R_p=R_L=\SI{10}{\kilo\ohm}$, tem-se
$xR_p\parallel R_L=\dfrac{10x}{1+x}\,\si{\kilo\ohm}$ e
\[
U_o=10\cdot\frac{\dfrac{10x}{1+x}}{10(1-x)+\dfrac{10x}{1+x}}
=10\cdot\frac{x}{(1-x)(1+x)+x}
=\frac{10x}{1+x-x^{2}}.
\]
(b)
\begin{center}
\begin{tabular}{cccc}
\hline
$x$ & ideal & com carga & desvio\\
\hline
$0{,}25$ & $\SI{2.50}{\volt}$ & $\SI{2.105}{\volt}$ & $-15{,}8\%$\\
$0{,}50$ & $\SI{5.00}{\volt}$ & $\SI{4.000}{\volt}$ & $-20{,}0\%$\\
$0{,}75$ & $\SI{7.50}{\volt}$ & $\SI{6.316}{\volt}$ & $-15{,}8\%$\\
\hline
\end{tabular}
\end{center}
(c) O desvio \emph{absoluto} é máximo perto do meio do curso, onde a
resistência de Thévenin do potenciômetro ($x(1-x)R_p$) atinge seu máximo,
$R_p/4=\SI{2.5}{\kilo\ohm}$. Nos extremos ($x\to0$ ou $x\to1$) a impedância
tende a zero e a carga deixa de importar.\\
\textbf{A curva deixa de ser reta:} o potenciômetro linear, carregado, comporta-se
de forma \textbf{não linear}, e o erro depende da posição --- o que é
inaceitável em um controle graduado.\\
\textbf{Correções:} \emph{(i)} usar $R_L\ge100R_p$, o que reduz o desvio máximo
a menos de $0{,}25\%$; \emph{(ii)} inserir um seguidor de tensão entre o cursor
e a carga, solução padrão em instrumentação; \emph{(iii)} usar potenciômetro de
lei "logarítmica reversa" pré-distorcida --- gambiarra que só funciona para uma
carga específica.}
\end{questao}

\begin{questao}[4]{}
\textbf{(Tolerância de razão.)} Um divisor usa resistores de $1\%$ com entrada
exata.
\begin{itensq}
\item Deduza a expressão do pior erro relativo da saída.
\item Avalie para as razões $12\to\SI{6}{\volt}$, $12\to\SI{5}{\volt}$ e
      $100\to\SI{5}{\volt}$.
\item Explique por que resistores \emph{casados} melhoram drasticamente o
      resultado.
\end{itensq}
\resposta{$\varepsilon=\dfrac{R_1}{R_1+R_2}(t_1+t_2)$; $1{,}00\%$,
$1{,}17\%$ e $1{,}90\%$}
\resolucao{(a) De $U_o=U_{in}\dfrac{R_2}{R_1+R_2}$, tomando logaritmos e
diferenciando:
\[
\frac{\Delta U_o}{U_o}
=\frac{R_1}{R_1+R_2}\left(\frac{\Delta R_2}{R_2}-\frac{\Delta R_1}{R_1}\right),
\]
cujo módulo máximo (desvios opostos) é
$\varepsilon=\dfrac{R_1}{R_1+R_2}(t_1+t_2)$.\\
(b) O fator $\dfrac{R_1}{R_1+R_2}$ vale $1-U_o/U_{in}$:
\begin{center}
\begin{tabular}{ccc}
\hline
Razão & $R_1/(R_1+R_2)$ & Pior erro\\
\hline
$12\to6$ & $0{,}500$ & $1{,}00\%$\\
$12\to5$ & $0{,}583$ & $1{,}17\%$\\
$100\to5$ & $0{,}950$ & $1{,}90\%$\\
\hline
\end{tabular}
\end{center}
\textbf{Padrão:} quanto \emph{maior} a razão de divisão, pior o erro --- ele
tende a $t_1+t_2=2\%$ quando $U_o\ll U_{in}$, e nunca o ultrapassa.\\
(c) \textbf{Resistores casados.} Em uma rede resistiva integrada, os dois
resistores são fabricados lado a lado, no mesmo processo e no mesmo substrato.
Seus desvios são \textbf{correlacionados}: se um sai $0{,}8\%$ acima do
nominal, o outro também sai. O termo entre parênteses,
$\dfrac{\Delta R_2}{R_2}-\dfrac{\Delta R_1}{R_1}$, tende a \emph{zero}.\\
Na prática, redes casadas oferecem razão com precisão de $0{,}05\%$ mesmo com
valores absolutos de $\pm20\%$. \textbf{O que importa no divisor é a razão, não
os valores} --- e o casamento é a única forma barata de obtê-la com precisão.
Esse é o princípio que viabiliza conversores A/D e amplificadores de
instrumentação.}
\end{questao}

\begin{questao}[4]{}
\textbf{(Sonda atenuadora.)} Uma sonda $10{:}1$ é formada por
$\SI{9}{\mega\ohm}$ na ponta, em série com a entrada de
$\SI{1}{\mega\ohm}$ do osciloscópio.
\begin{itensq}
\item Verifique a razão de atenuação e a resistência de entrada total.
\item Determine o erro ao medir uma fonte com $R_{Th}=\SI{100}{\kilo\ohm}$,
      com e sem a sonda.
\item Explique o que se perde ao usar a sonda.
\end{itensq}
\resposta{(a) $1/10$ e $\SI{10}{\mega\ohm}$; (b) $1{,}0\%$ com sonda,
$9{,}1\%$ sem}
\resolucao{(a) A entrada do osciloscópio \emph{é} o $R_2$ do divisor:
\[
\frac{U_o}{U_{in}}=\frac{1}{9+1}=\frac{1}{10}\ \checkmark;
\qquad
R_{entrada}=9+1=\SI{10}{\mega\ohm}.
\]
(b) \textbf{Com sonda:} o instrumento apresenta $\SI{10}{\mega\ohm}$ à fonte:
\[
\varepsilon=\frac{100\,\mathrm{k}}{10\,\mathrm{M}+100\,\mathrm{k}}=1{,}0\%.
\]
\textbf{Sem sonda} (cabo direto, $\SI{1}{\mega\ohm}$):
\[
\varepsilon=\frac{100\,\mathrm{k}}{1\,\mathrm{M}+100\,\mathrm{k}}=9{,}1\%.
\]
A sonda reduz o erro de carregamento por um fator $9$.\\
(c) \textbf{O que se perde: amplitude.} O sinal chega ao instrumento dez vezes
menor, o que reduz a relação sinal-ruído em $\SI{20}{\decibel}$ e torna
inviável medir sinais de poucos milivolts.\\
\textbf{O compromisso central da sonda} é exatamente esse: troca-se sensibilidade
por não perturbação. Por isso os osciloscópios trazem seletor $1\times/10\times$
--- sinais pequenos em fontes de baixa impedância pedem $1\times$; sinais grandes
em circuitos de alta impedância pedem $10\times$.\\
\emph{(Em sinais variáveis há ainda a capacitância parasita, que exige o
capacitor de compensação ajustável presente em toda sonda real --- assunto de
corrente alternada.)}}
\end{questao}

\begin{questao}[4]{}
\textbf{(Otimização de sensor.)} Um termistor de resistência $R_t$ ocupa a
posição inferior de um divisor com resistor fixo $R$ e alimentação $U$.
\begin{itensq}
\item Obtenha $\dfrac{\mathrm{d}U_o}{\mathrm{d}R_t}$.
\item Determine o valor de $R$ que maximiza a sensibilidade em um dado ponto de
      operação.
\item Avalie para $U=\SI{5}{\volt}$, $R_t=\SI{10}{\kilo\ohm}$ e coeficiente de
      $\SI{-400}{\ohm\per\celsius}$.
\end{itensq}
\resposta{(b) $R=R_t$; (c) $\SI{-50}{\milli\volt\per\celsius}$}
\resolucao{(a) De $U_o=U\dfrac{R_t}{R+R_t}$:
\[
\frac{\mathrm{d}U_o}{\mathrm{d}R_t}
=U\,\frac{(R+R_t)-R_t}{(R+R_t)^{2}}
=\frac{U\,R}{(R+R_t)^{2}}.
\]
(b) Maximizando em relação a $R$, com $R_t$ fixo:
\[
\frac{\mathrm{d}}{\mathrm{d}R}\left[\frac{R}{(R+R_t)^{2}}\right]
=\frac{(R+R_t)-2R}{(R+R_t)^{3}}
=\frac{R_t-R}{(R+R_t)^{3}}=0
\;\Longrightarrow\;
\boxed{R=R_t}
\]
--- e a saída no ponto ótimo é exatamente $U/2$.\\
\textbf{Mesma estrutura} do resultado de máxima potência: a derivada de
$R/(R+R_t)^{2}$ é a mesma que a de $P_R$ em uma série, e por isso o ótimo cai
em igualdade.\\
(c) Com $R=R_t=\SI{10}{\kilo\ohm}$:
\[
\frac{\mathrm{d}U_o}{\mathrm{d}R_t}=\frac{5(10^{4})}{(\pot{2}{4})^{2}}
=\SI{1.25e-4}{\volt\per\ohm},
\]
\[
\frac{\mathrm{d}U_o}{\mathrm{d}T}
=(\pot{1.25}{-4})(-400)=\SI{-50}{\milli\volt\per\celsius}.
\]
\textbf{Excelente sensibilidade:} um conversor A/D de $10$ bits com fundo de
$\SI{5}{\volt}$ tem degrau de $\SI{4.9}{\milli\volt}$, resolvendo cerca de
$\SI{0.1}{\celsius}$.\\
\textbf{Ressalva:} o ótimo vale para \emph{um} ponto de operação. Em faixa
ampla de temperatura, $R_t$ varia por ordens de grandeza e a sensibilidade cai
nas extremidades --- escolha $R$ igual a $R_t$ na temperatura de maior interesse,
não na média da faixa.}
\end{questao}

\begin{questao}[4]{}
\textbf{(Divisor contra regulador.)} Compare as duas soluções para obter
$\SI{5}{\volt}$ a partir de $\SI{12}{\volt}$, para uma carga que varia de
$\SI{1}{\milli\ampere}$ a $\SI{10}{\milli\ampere}$.
\begin{itensq}
\item Dimensione o divisor com repouso de dez vezes a carga máxima e determine
      a regulação.
\item Determine o rendimento de cada solução.
\item Estabeleça o critério de escolha.
\end{itensq}
\resposta{(a) $R_1=\SI{70}{\ohm}$, $R_2=\SI{50}{\ohm}$; regulação
$\approx5{,}3\%$; (b) $3{,}8\%$ contra $42\%$}
\resolucao{(a) Repouso de $\SI{100}{\milli\ampere}$:
$R_2=5/0{,}1=\SI{50}{\ohm}$ e
$R_1=7/0{,}1=\SI{70}{\ohm}$.
\[
R_{Th}=50\parallel70=\SI{29.2}{\ohm};
\qquad
\Delta U=R_{Th}\,\Delta I_L=29{,}2(\pot{9}{-3})=\SI{0.26}{\volt},
\]
ou seja, cerca de $5{,}3\%$ de variação.\\
(b) \textbf{Divisor:} consumo total
$\approx12\times0{,}11=\SI{1.32}{\watt}$ para entregar no máximo
$5\times0{,}01=\SI{50}{\milli\watt}$ --- rendimento de $3{,}8\%$.\\
\textbf{Regulador linear:} consome a mesma corrente da carga,
$12\times0{,}01=\SI{120}{\milli\watt}$ para entregar
$\SI{50}{\milli\watt}$ --- rendimento de $41{,}7\%$ (igual à razão
$U_o/U_{in}$), com regulação típica melhor que $0{,}1\%$.\\
(c) \textbf{Critério.}
\begin{itemize}
\item \textbf{Divisor} apenas quando a carga é de \emph{altíssima impedância} e
praticamente constante: entradas de conversores A/D, polarização de bases,
referências internas. Vantagens: dois componentes, custo desprezível, sem ruído
de comutação.
\item \textbf{Regulador} sempre que a carga consumir corrente apreciável ou
variar. O regulador linear resolve com rendimento $U_o/U_{in}$; o chaveado
ultrapassa $90\%$, ao custo de complexidade e ruído.
\end{itemize}
\textbf{Fronteira prática:} se a corrente da carga for maior que cerca de
$1\%$ da corrente de repouso que o divisor exigiria, o divisor já não é a
escolha certa.}
\end{questao}

\begin{questao}[4]{}
\textbf{(Divisor com proteção.)} Projete a interface entre um sinal de entrada
que varia de $\SI{0}{\volt}$ a $\SI{30}{\volt}$ e um conversor A/D de
$\SI{3.3}{\volt}$ e impedância de entrada muito alta.
\begin{itensq}
\item Determine a razão e escolha os resistores.
\item Verifique o comportamento em sobretensão de $\SI{50}{\volt}$ e proponha
      proteção.
\item Determine o consumo e discuta o compromisso com a impedância exigida pelo
      conversor.
\end{itensq}
\resposta{(a) razão $\approx1/10$: $R_1=\SI{91}{\kilo\ohm}$,
$R_2=\SI{10}{\kilo\ohm}$; (c) $\SI{297}{\micro\ampere}$ no fundo de escala}
\resolucao{(a) É preciso $30\to\SI{3.3}{\volt}$, razão $0{,}11$. Com
$R_2=\SI{10}{\kilo\ohm}$ e $R_1=\SI{91}{\kilo\ohm}$ (ambos E24):
\[
U_o=30\cdot\frac{10}{101}=\SI{2.97}{\volt},
\]
$10\%$ abaixo do fundo de escala --- margem prudente, que acomoda tolerância dos
resistores sem saturar o conversor.\\
(b) Com $\SI{50}{\volt}$ na entrada, a saída iria a
$50\times10/101=\SI{4.95}{\volt}$, acima do limite absoluto do conversor.\\
\textbf{Proteção:} um diodo Zener de $\SI{3.3}{\volt}$ (ou um diodo de
grampeamento para a alimentação) em paralelo com $R_2$. Em sobretensão, ele
conduz e limita a saída; a corrente é confinada por $R_1$ a
\[
I=\frac{50-3{,}3}{91\,\mathrm{k}}=\SI{0.51}{\milli\ampere},
\]
valor inofensivo tanto para o Zener quanto para $R_1$
($P=\SI{24}{\milli\watt}$). \textbf{É $R_1$ que faz a proteção funcionar} ---
sem ele, o diodo teria de absorver corrente ilimitada.\\
(c) No fundo de escala, $I=30/101\,\mathrm{k}=\SI{297}{\micro\ampere}$ e
$P\approx\SI{8.9}{\milli\watt}$.\\
\textbf{Compromisso.} $R_{Th}=91\parallel10=\SI{9}{\kilo\ohm}$. Conversores A/D
de aproximações sucessivas exigem impedância de fonte baixa (tipicamente
$<\SI{10}{\kilo\ohm}$) para carregar o capacitor de amostragem no tempo
disponível --- este projeto fica no limite.\\
\textbf{Soluções:} reduzir os resistores por um fator $10$ (custa dez vezes mais
consumo), acrescentar um capacitor de $\SI{100}{\nano\farad}$ na entrada do
conversor (fornece a carga instantânea de amostragem), ou inserir um seguidor
de tensão. A segunda opção é a mais comum, e a mais barata.}
\end{questao}

% END BANCO COMPLEMENTAR

\gabarito[4.1 Divisor de tensão]

% =====================================================================
\subsection{Divisor de corrente}
% =====================================================================

\blocofixacao

\begin{questao}[1]{}
Em um divisor de corrente de dois resistores em paralelo, a maior corrente
circula:
\begin{alternativas}
\item pelo maior resistor.
\item pelo menor resistor.
\item igualmente pelos dois, sempre.
\item pelo resistor mais próximo da fonte.
\item por nenhum: toda a corrente retorna à fonte.
\end{alternativas}
\resposta{(b)}
\resolucao{Ambos os ramos têm a mesma tensão; por $I = U/R$, o menor resistor
conduz a maior corrente. A corrente se reparte de forma \emph{inversamente}
proporcional às resistências.}
\end{questao}

\begin{questao}[1]{}
Uma corrente de $\SI{10}{\ampere}$ chega a um nó e se divide entre
$R_1=\SI{2}{\ohm}$ e $R_2=\SI{3}{\ohm}$ em paralelo. Calcule $i_1$ e $i_2$.

\circuitoq{%
\begin{circuitikz}[scale=1,transform shape,line width=0.8pt,american]
 \draw (0,0) to[isource,l={$\SI{10}{\ampere}$}] (0,2.6)
       to[short,-*,i={$I_t$}] (2,2.6)
       to[R,l={$R_1=\SI{2}{\ohm}$},i>_={$i_1$}] (2,0) -- (0,0);
 \draw (2,2.6) -- (4,2.6)
       to[R,l={$R_2=\SI{3}{\ohm}$},i>_={$i_2$}] (4,0) to[short,-*] (2,0);
\end{circuitikz}}

\resposta{$i_1 = \SI{6}{\ampere}$; $i_2 = \SI{4}{\ampere}$}
\resolucao{Pelo divisor de corrente (índices \emph{trocados}):
\[
i_1 = I_t\,\frac{R_2}{R_1+R_2}=10\cdot\frac{3}{5}=\SI{6}{\ampere},
\qquad
i_2 = I_t\,\frac{R_1}{R_1+R_2}=10\cdot\frac{2}{5}=\SI{4}{\ampere}.
\]
\textbf{Verificação:} $6+4 = \SI{10}{\ampere}$; e o menor resistor
(\SI{2}{\ohm}) ficou com a maior corrente.}
\end{questao}

\begin{questao}[1]{}
Uma corrente de $\SI{15}{\ampere}$ divide-se entre dois resistores em paralelo,
$\SI{3}{\ohm}$ e $\SI{6}{\ohm}$. As correntes valem, respectivamente:
\begin{alternativas}[2]
\item \SI{5}{\ampere} e \SI{10}{\ampere}
\item \SI{10}{\ampere} e \SI{5}{\ampere}
\item \SI{7.5}{\ampere} e \SI{7.5}{\ampere}
\item \SI{9}{\ampere} e \SI{6}{\ampere}
\item \SI{6}{\ampere} e \SI{9}{\ampere}
\end{alternativas}
\resposta{(b)}
\resolucao{$i_{3\Omega} = 15\cdot\dfrac{6}{9} = \SI{10}{\ampere}$ e
$i_{6\Omega} = 15\cdot\dfrac{3}{9} = \SI{5}{\ampere}$. O resistor de \SI{3}{\ohm},
por ser menor, conduz o dobro do de \SI{6}{\ohm}.}
\end{questao}

\blocoaplicacao

\begin{questao}[2]{}
No circuito, uma fonte de corrente de $\SI{30}{\ampere}$ alimenta três resistores
em paralelo. Calcule a corrente em cada um.

\circuitoq{%
\begin{circuitikz}[scale=1,transform shape,line width=0.8pt,american]
 \draw (0,0) to[isource,l={$\SI{30}{\ampere}$}] (0,2.6)
       to[short,-*,i={$I_t$}] (2,2.6)
       to[R,l={$\SI{2}{\ohm}$},i>_={$i_1$}] (2,0) -- (0,0);
 \draw (2,2.6) -- (4,2.6)
       to[R,l={$\SI{3}{\ohm}$},i>_={$i_2$}] (4,0) to[short,-*] (2,0);
 \draw (4,2.6) -- (6,2.6)
       to[R,l={$\SI{6}{\ohm}$},i>_={$i_3$}] (6,0) to[short,-*] (4,0);
\end{circuitikz}}

\resposta{$i_1=\SI{15}{\ampere}$, $i_2=\SI{10}{\ampere}$, $i_3=\SI{5}{\ampere}$}
\resolucao{Com três ou mais ramos, o mais prático é achar a tensão comum.
$\Req = 2\parallel3\parallel6$: como
$\dfrac{1}{2}+\dfrac{1}{3}+\dfrac{1}{6}=\dfrac{3+2+1}{6}=1$, temos
$\Req=\SI{1}{\ohm}$ e $U = 30\cdot1 = \SI{30}{\volt}$. Então:
\[
i_1=\frac{30}{2}=\SI{15}{\ampere},\quad
i_2=\frac{30}{3}=\SI{10}{\ampere},\quad
i_3=\frac{30}{6}=\SI{5}{\ampere}.
\]
\textbf{Verificação:} $15+10+5 = \SI{30}{\ampere}$.}
\end{questao}

\begin{questao}[2]{}
No divisor da figura, com três ramos, determine a corrente em cada resistor.

\circuitoq{%
\begin{circuitikz}[scale=1,transform shape,line width=0.8pt,american]
 \draw (0,0) to[isource,l={$\SI{25}{\ampere}$}] (0,2.6)
       to[short,-*,i={$I_t$}] (2,2.6)
       to[R,l={$\SI{12}{\ohm}$},i>_={$i_1$}] (2,0) -- (0,0);
 \draw (2,2.6) -- (4,2.6)
       to[R,l={$\SI{3}{\ohm}$},i>_={$i_2$}] (4,0) to[short,-*] (2,0);
 \draw (4,2.6) -- (6,2.6)
       to[R,l={$\SI{12}{\ohm}$},i>_={$i_3$}] (6,0) to[short,-*] (4,0);
\end{circuitikz}}

\resposta{$i_1=i_3\approx\SI{4.17}{\ampere}$; $i_2\approx\SI{16.7}{\ampere}$}
\resolucao{$\dfrac{1}{\Req}=\dfrac{1}{12}+\dfrac{1}{3}+\dfrac{1}{12}
=\dfrac{1+4+1}{12}=\dfrac{6}{12}$, logo $\Req=\SI{2}{\ohm}$ e
$U = 25\cdot2 = \SI{50}{\volt}$.
\[
i_1=i_3=\frac{50}{12}\approx\SI{4.17}{\ampere},
\qquad
i_2=\frac{50}{3}\approx\SI{16.7}{\ampere}.
\]
Os dois ramos de \SI{12}{\ohm} conduzem o mesmo (por simetria); o ramo de
\SI{3}{\ohm}, quatro vezes menor, conduz quatro vezes mais.}
\end{questao}

\blocoaprofundamento

\begin{questao}[3]{}
Uma corrente de $\SI{35}{\ampere}$ alimenta três resistores em paralelo,
$\SI{5}{\ohm}$, $\SI{10}{\ohm}$ e $\SI{20}{\ohm}$.
\begin{itensq}
\item Calcule a corrente em cada ramo.
\item Verifique que a soma das potências dissipadas nos ramos é igual à potência
      entregue pela fonte.
\end{itensq}
\resposta{(a) \SI{20}{\ampere}, \SI{10}{\ampere}, \SI{5}{\ampere};
(b) \SI{3500}{\watt} nos dois lados}
\resolucao{(a) $\dfrac{1}{\Req}=\dfrac{1}{5}+\dfrac{1}{10}+\dfrac{1}{20}
=\dfrac{4+2+1}{20}=\dfrac{7}{20}$, então $\Req = \dfrac{20}{7}\,\si{\ohm}$ e
$U = 35\cdot\dfrac{20}{7} = \SI{100}{\volt}$. As correntes:
\[
i_5=\frac{100}{5}=\SI{20}{\ampere},\quad
i_{10}=\frac{100}{10}=\SI{10}{\ampere},\quad
i_{20}=\frac{100}{20}=\SI{5}{\ampere}.
\]
(b) Potência em cada ramo ($P = U\,i$, com $U=\SI{100}{\volt}$):
\[
100\cdot20 + 100\cdot10 + 100\cdot5 = 2000+1000+500 = \SI{3500}{\watt}.
\]
Potência da fonte: $P = U\,I_t = 100\cdot35 = \SI{3500}{\watt}$ --- confere com a
conservação de energia.}
\end{questao}

\begin{questao}[3]{}
Um miliamperímetro tem fundo de escala de $\SI{1}{\milli\ampere}$ e resistência
interna $R_g = \SI{50}{\ohm}$. Para medir correntes de até $\SI{1}{\ampere}$,
liga-se em paralelo com ele um resistor \emph{shunt} $R_s$, que desvia o excesso
de corrente. Determine $R_s$.

\circuitoq{%
\begin{circuitikz}[scale=1,transform shape,line width=0.8pt,american]
 \draw (0,1.4) to[short,i={$I=\SI{1}{\ampere}$},-*] (2,1.4) coordinate(n1);
 \draw (n1) to[R,l={$R_g=\SI{50}{\ohm}$},i={$\SI{1}{\milli\ampere}$}] (5,1.4) coordinate(n2);
 \node[above,font=\footnotesize] at (3.5,1.75) {galvanômetro};
 \draw (n1) to[short] (2,0) to[R,l_={$R_s$},i_={$i_s$}] (5,0) to[short,-*] (n2);
 \draw (n2) to[short] (7,1.4);
\end{circuitikz}}

\resposta{$R_s \approx \SI{0.05}{\ohm}$}
\resolucao{O galvanômetro e o shunt estão em paralelo, logo têm a \emph{mesma
tensão}. No fundo de escala, o medidor conduz $\SI{1}{\milli\ampere}$ e o shunt
deve desviar o restante:
\[
i_s = 1 - 0{,}001 = \SI{0.999}{\ampere}.
\]
Igualando as tensões nos dois ramos ($R_g\,i_g = R_s\,i_s$):
\[
R_s = \frac{R_g\,i_g}{i_s} = \frac{50\cdot0{,}001}{0{,}999}\approx\SI{0.05}{\ohm}.
\]
\textbf{Ideia central:} o shunt é um divisor de corrente que amplia a faixa do
instrumento por um fator $\dfrac{I}{i_g} = \dfrac{1}{0{,}001} = 1000$. É assim que
um mesmo galvanômetro vira amperímetro de várias escalas.}
\end{questao}

\begin{questao}[3]{}
Um galvanômetro tem fundo de escala $i_g = \SI{1}{\milli\ampere}$ e resistência
interna $R_g = \SI{50}{\ohm}$. Deseja-se convertê-lo em um amperímetro de
\textbf{três escalas}: \SI{100}{\milli\ampere}, \SI{1}{\ampere} e
\SI{10}{\ampere}, usando shunts comutáveis.
\begin{itensq}
\item Deduza a expressão geral do shunt para uma escala $I_{\max}$.
\item Calcule os três shunts.
\item Explique por que a resistência do amperímetro \emph{diminui} conforme se
      aumenta a escala, e por que isso é desejável.
\end{itensq}
\resposta{(a) $R_s = \dfrac{R_g\,i_g}{I_{\max}-i_g}$;
(b) \SI{0.505}{\ohm}, \SI{50.05}{\milli\ohm}, \SI{5.0}{\milli\ohm}; (c) ver comentário}
\resolucao{\textbf{(a)} O shunt está em paralelo com o galvanômetro, logo ambos
têm a mesma tensão. No fundo de escala, o galvanômetro conduz $i_g$ e o shunt
desvia o restante, $I_{\max}-i_g$:
\[
R_g\,i_g = R_s\,(I_{\max}-i_g)
\;\Longrightarrow\;
\boxed{\,R_s=\frac{R_g\,i_g}{I_{\max}-i_g}\,}
\]
\textbf{(b)} Aplicando com $R_g i_g = 50\cdot\pot{1}{-3}=\SI{0.05}{\volt}$:
\begin{center}\small\sffamily
\begin{tabular}{@{}c c c@{}}
\toprule
\textbf{Escala} & \textbf{$R_s$} & \textbf{Fator de multiplicação}\\
\midrule
\SI{100}{\milli\ampere} & $\dfrac{0{,}05}{0{,}099}\approx\SI{0.505}{\ohm}$ & 100$\times$\\
\SI{1}{\ampere}   & $\dfrac{0{,}05}{0{,}999}\approx\SI{50.05}{\milli\ohm}$ & 1000$\times$\\
\SI{10}{\ampere}  & $\dfrac{0{,}05}{9{,}999}\approx\SI{5.0}{\milli\ohm}$ & 10\,000$\times$\\
\bottomrule
\end{tabular}
\end{center}
\textbf{(c)} A resistência total do instrumento é $R_g \parallel R_s$. Como $R_s$
diminui a cada escala maior, o paralelo também diminui --- na escala de
\SI{10}{\ampere}, o amperímetro apresenta apenas $\approx\SI{5}{\milli\ohm}$.\\
\textbf{Por que isso é desejável:} o amperímetro é ligado \emph{em série} com o
circuito, então sua resistência se soma à do circuito e \emph{altera} a corrente
que se pretende medir --- o chamado \textbf{erro de inserção}. Quanto menor a
resistência do instrumento, menor a perturbação. O amperímetro ideal teria
resistência nula (e o voltímetro ideal, resistência infinita).\\
\textbf{Estimativa do erro:} medindo \SI{10}{\ampere} em um circuito de
$\SI{1}{\ohm}$, o erro relativo seria $\dfrac{0{,}005}{1}=\SI{0.5}{\percent}$ ---
aceitável. Mas em um circuito de $\SI{0.05}{\ohm}$, o erro subiria para
\SI{10}{\percent}, tornando a medida inútil.}
\end{questao}

% BEGIN BANCO COMPLEMENTAR Divisor de corrente
% =====================================================================
%  Banco complementar --- Divisor de Corrente  (REVISADO)
%  Substitui a versao gerada por template (10 clones em N2, 9 em N3 e
%  8 questoes conceituais marcadas como N4).
%  RESTRICAO FORTE: o projeto de SHUNT ja aparece DUAS vezes no cap04
%  (miliamperimetro 1 mA/50 ohm e galvanometro 1 mA/50 ohm) e mais
%  duas vezes no banco de circuitos paralelo revisado (shunt 99 ohm e
%  shunt de Ayrton multiescala). Este banco NAO trata de shunt.
%  O cap04 tambem cobre: maior corrente (MC), 10 A entre 2/3, 15 A
%  entre 3/6, 30 A em tres paralelos, tres ramos com figura e 35 A
%  entre 5/10/20 -- de modo que a divisao basica de dois e tres ramos
%  tambem esta excluida.
%  Distribuicao mantida: 5 N1 + 10 N2 + 9 N3 + 8 N4 = 32
% =====================================================================

\subsubsection*{Banco complementar --- Divisor de corrente}

\begin{atencao}[Organização do banco]
A forma correta e geral do divisor de corrente é em \textbf{condutâncias}:
$I_k=I_T\dfrac{G_k}{\sum G_j}$. A versão com resistências,
$I_1=I_T\dfrac{R_2}{R_1+R_2}$, é apenas o caso particular de \emph{dois} ramos
--- e não se generaliza. O N3 trata de desequilíbrio entre condutores,
tolerância, falhas e o contraste entre alimentação por fonte de corrente e por
fonte de tensão; o N4 exige demonstração, projeto de balanceamento e análise de
sensibilidade.
\end{atencao}

\blocofixacao

\begin{questao}[1]{}
Três ramos com condutâncias $G_1=\SI{0.5}{\siemens}$,
$G_2=\SI{0.3}{\siemens}$ e $G_3=\SI{0.2}{\siemens}$ recebem
$I_T=\SI{10}{\ampere}$. Determine a corrente de cada ramo.
\resposta{$\SI{5}{\ampere}$, $\SI{3}{\ampere}$ e $\SI{2}{\ampere}$}
\resolucao{Com $\sum G=0{,}5+0{,}3+0{,}2=\SI{1}{\siemens}$:
\[
I_k=I_T\frac{G_k}{\sum G}
\;\Longrightarrow\;
5;\ 3;\ \SI{2}{\ampere}.
\]
Em condutâncias a conta é imediata: a corrente se reparte na \emph{mesma
proporção} das condutâncias. É por isso que essa é a forma natural do divisor de
corrente, e não a expressão com resistências.}
\end{questao}

\begin{questao}[1]{}
Dois ramos de $\SI{4}{\ohm}$ e $\SI{12}{\ohm}$ dividem uma corrente. Determine
a razão $I_1/I_2$.
\resposta{$I_1/I_2=3$}
\resolucao{Como a tensão é comum, $I=U/R$ e
\[
\frac{I_1}{I_2}=\frac{R_2}{R_1}=\frac{12}{4}=3.
\]
A razão das correntes é a \textbf{inversa} da razão das resistências --- o ramo
de menor resistência leva a maior corrente. Compare com a série, em que a razão
das \emph{tensões} é \emph{direta}.}
\end{questao}

\begin{questao}[1]{}
Um ramo de $\SI{20}{\ohm}$ de uma associação paralela conduz
$\SI{1.5}{\ampere}$; o outro ramo vale $\SI{30}{\ohm}$. Determine a corrente
total.
\resposta{$I_T=\SI{2.5}{\ampere}$}
\resolucao{A tensão comum sai do ramo conhecido:
\[
U=20(1{,}5)=\SI{30}{\volt}
\;\Longrightarrow\;
I_2=\frac{30}{30}=\SI{1}{\ampere},
\]
\[
I_T=1{,}5+1=\SI{2.5}{\ampere}.
\]
Um único ramo conhecido determina a tensão comum e, com ela, todo o resto ---
propriedade que torna qualquer resistor de um paralelo um sensor de corrente em
potencial.}
\end{questao}

\begin{questao}[1]{}
Um dos ramos de um divisor de corrente sofre \textbf{curto-circuito}. Determine
o que ocorre com as correntes.
\resposta{Toda a corrente passa pelo curto; os demais ramos ficam sem corrente}
\resolucao{O ramo em curto tem $R=0$, logo $G\to\infty$. Pelo divisor,
\[
I_k=I_T\frac{G_k}{\infty+\sum_{j\ne k}G_j}\to0
\]
para todos os ramos finitos, e $I_{curto}\to I_T$.\\
\textbf{Fisicamente:} o curto impõe $U=0$ à associação inteira, e sem tensão
nenhum resistor conduz. Os demais ramos continuam íntegros --- eles apenas
deixam de funcionar, sendo vítimas e não causa do defeito.}
\end{questao}

\begin{questao}[1]{}
Em um divisor de corrente alimentado por uma \textbf{fonte de corrente ideal},
um dos ramos é aberto. A corrente nos ramos restantes:
\begin{alternativas}[2]
\item permanece a mesma
\item aumenta
\item diminui
\item torna-se nula
\end{alternativas}
\resposta{(b)}
\resolucao{A fonte de corrente ideal mantém $I_T$ fixo. Com um caminho a menos,
a mesma corrente se distribui entre menos ramos --- cada um recebe \emph{mais}.\\
\textbf{Contraste essencial:} sob fonte de \emph{tensão} ideal, os ramos
restantes ficariam \emph{inalterados} (a tensão comum não muda) e apenas a
corrente total cairia. A mesma falha produz efeitos opostos conforme o tipo de
alimentação --- ponto explorado quantitativamente no bloco de
aprofundamento.}
\end{questao}

\blocoaplicacao

\begin{questao}[2]{}
Uma corrente de $\SI{12}{\ampere}$ divide-se entre quatro ramos de
$\SI{5}{\ohm}$, $\SI{10}{\ohm}$, $\SI{20}{\ohm}$ e $\SI{20}{\ohm}$.
\begin{itensq}
\item Determine a tensão da associação.
\item Determine a corrente de cada ramo.
\end{itensq}
\resposta{$U=\SI{30}{\volt}$; $6$, $3$, $1{,}5$ e $\SI{1.5}{\ampere}$}
\resolucao{
\[
\sum G=0{,}20+0{,}10+0{,}05+0{,}05=\SI{0.40}{\siemens}
\;\Longrightarrow\;
U=\frac{I_T}{\sum G}=\frac{12}{0{,}40}=\SI{30}{\volt}.
\]
\[
I_k=\frac{U}{R_k}\Rightarrow 6;\ 3;\ 1{,}5;\ \SI{1.5}{\ampere}.
\]
\textbf{Verificação:} $6+3+1{,}5+1{,}5=\SI{12}{\ampere}$ \checkmark.\\
\textbf{Método:} calcular a tensão comum primeiro é quase sempre mais rápido do
que aplicar a fórmula do divisor ramo a ramo --- e produz um valor intermediário
que serve de conferência.}
\end{questao}

\begin{questao}[2]{}
No circuito anterior, determine a fração percentual da corrente em cada ramo e
explique por que ela não depende de $I_T$.
\resposta{$50\%$, $25\%$, $12{,}5\%$ e $12{,}5\%$}
\resolucao{
\[
\frac{I_k}{I_T}=\frac{G_k}{\sum G}
\Rightarrow \frac{0{,}20}{0{,}40}=50\%;\quad25\%;\quad12{,}5\%;\quad12{,}5\%.
\]
A razão depende \emph{apenas} das condutâncias --- $I_T$ cancela. Dobrar a
corrente total dobra todas as parcelas, preservando as proporções.\\
\textbf{Consequência prática:} o perfil de divisão é uma propriedade fixa da
rede. Ele só muda se algum ramo mudar de valor, for aberto ou entrar em curto
--- o que faz da monitoração dessas frações uma técnica de diagnóstico.}
\end{questao}

\begin{questao}[2]{}
Uma corrente de $\SI{8}{\ampere}$ deve dividir-se de modo que
$\SI{25}{\percent}$ passe por um ramo $R$ e o restante por um ramo de
$\SI{15}{\ohm}$. Determine $R$.
\resposta{$R=\SI{45}{\ohm}$}
\resolucao{As correntes são $I_R=\SI{2}{\ampere}$ e
$I_{15}=\SI{6}{\ampere}$. A tensão é comum:
\[
U=15(6)=\SI{90}{\volt}
\;\Longrightarrow\;
R=\frac{90}{2}=\SI{45}{\ohm}.
\]
\textbf{Conferência pela razão:} $I_R/I_{15}=R_{15}/R=15/45=1/3$, e de fato
$2/6=1/3$ \checkmark.\\
Note que $R$ é o \emph{triplo} do outro ramo para conduzir um \emph{terço} da
corrente dele --- a proporcionalidade inversa em ação.}
\end{questao}

\begin{questao}[2]{}
Uma fonte de $\SI{60}{\volt}$ com resistência série $R_s=\SI{4}{\ohm}$ alimenta
dois ramos em paralelo de $\SI{10}{\ohm}$ e $\SI{15}{\ohm}$.
\begin{itensq}
\item Determine a corrente total.
\item Determine a corrente de cada ramo.
\end{itensq}
\resposta{$I_T=\SI{6}{\ampere}$; $\SI{3.6}{\ampere}$ e $\SI{2.4}{\ampere}$}
\resolucao{(a) $R_{par}=\dfrac{10\cdot15}{25}=\SI{6}{\ohm}$ e
\[
I_T=\frac{60}{4+6}=\SI{6}{\ampere}.
\]
(b) Pelo divisor de dois ramos:
\[
I_1=6\cdot\frac{15}{25}=\SI{3.6}{\ampere};
\qquad
I_2=6\cdot\frac{10}{25}=\SI{2.4}{\ampere}.
\]
\textbf{Ponto de atenção:} $R_s$ afeta a corrente \emph{total}, mas não a
\emph{proporção} da divisão --- ela é determinada só pelos dois ramos. Essa
separação entre "quanto chega" e "como se reparte" simplifica muito a análise de
circuitos reais.}
\end{questao}

\begin{questao}[2]{}
No circuito anterior, determine a potência dissipada em cada ramo e a tensão
sobre a associação.
\resposta{$U=\SI{36}{\volt}$; $\SI{129.6}{\watt}$ e $\SI{86.4}{\watt}$}
\resolucao{
\[
U=R_{par}I_T=6(6)=\SI{36}{\volt};
\qquad
P_1=\frac{36^{2}}{10}=\SI{129.6}{\watt};
\qquad
P_2=\frac{36^{2}}{15}=\SI{86.4}{\watt}.
\]
Total nos ramos: $\SI{216}{\watt}=36(6)$ \checkmark.\\
\textbf{Observe:} o ramo de \emph{menor} resistência dissipa mais --- regra do
paralelo ($P\propto1/R$ com $U$ comum). O resistor série dissipa ainda
$4(36)=\SI{144}{\watt}$, e a fonte entrega $60(6)=\SI{360}{\watt}$
\checkmark.}
\end{questao}

\begin{questao}[2]{}
Uma corrente de $\SI{10}{\ampere}$ divide-se entre $\SI{1}{\ohm}$ e
$\SI{1}{\kilo\ohm}$.
\begin{itensq}
\item Determine as duas correntes.
\item Avalie o erro de simplesmente desprezar o ramo de alta resistência.
\end{itensq}
\resposta{$\SI{9.99}{\ampere}$ e $\SI{9.99}{\milli\ampere}$; erro de
$0{,}1\%$}
\resolucao{(a)
\[
I_1=10\cdot\frac{1000}{1001}=\SI{9.990}{\ampere};
\qquad
I_2=10\cdot\frac{1}{1001}=\SI{9.99}{\milli\ampere}.
\]
(b) Desprezar o segundo ramo daria $I_1=\SI{10}{\ampere}$ --- erro de
$0{,}1\%$, muito abaixo da tolerância de qualquer resistor real.\\
\textbf{Regra prática:} um ramo mil vezes maior desvia um milésimo da corrente.
Se a razão for de $100{:}1$, o erro sobe para $1\%$; se for $10{:}1$, para
$9\%$ --- aí já não se pode desprezar.\\
\textbf{Mas atenção:} desprezar o ramo para calcular $I_1$ é legítimo;
\emph{esquecê-lo} pode não ser, pois aqueles $\SI{10}{\milli\ampere}$ podem ser
exatamente o sinal que interessa medir.}
\end{questao}

\begin{questao}[2]{}
Uma fonte de corrente de $\SI{12}{\ampere}$ alimenta três ramos de
$\SI{6}{\ohm}$, $\SI{12}{\ohm}$ e $\SI{12}{\ohm}$. O ramo de
$\SI{6}{\ohm}$ é removido. Determine as novas correntes.
\resposta{Antes: $6$, $3$, $\SI{3}{\ampere}$. Depois: $\SI{6}{\ampere}$ em cada}
\resolucao{\textbf{Antes:} $\sum G=\frac16+\frac1{12}+\frac1{12}
=\SI{0.3333}{\siemens}$, $U=12/0{,}3333=\SI{36}{\volt}$ e
$I=6;\ 3;\ \SI{3}{\ampere}$.\\
\textbf{Depois:} $\sum G=\SI{0.1667}{\siemens}$,
$U=12/0{,}1667=\SI{72}{\volt}$ e $I=\SI{6}{\ampere}$ em cada.\\
\textbf{A tensão dobrou} e a corrente de cada ramo remanescente também. Note que
os resistores não mudaram --- foi a \emph{fonte de corrente} que, obrigada a
manter $\SI{12}{\ampere}$, elevou a tensão para forçar a mesma corrente por
menos caminhos.\\
\textbf{Alerta de projeto:} em circuitos alimentados por fonte de corrente,
abrir um ramo \emph{sobrecarrega} os demais. É o oposto do que ocorre com fonte
de tensão.}
\end{questao}

\begin{questao}[2]{}
Um resistor de $\SI{10}{\ohm}$ está em paralelo com um \textbf{indutor} ideal,
e o conjunto recebe $\SI{4}{\ampere}$ em regime permanente CC. Determine a
corrente de cada ramo.
\resposta{$\SI{0}{\ampere}$ no resistor; $\SI{4}{\ampere}$ no indutor}
\resolucao{Em regime permanente CC, a corrente do indutor é constante e
$v=L\,\mathrm{d}i/\mathrm{d}t=0$: ele se comporta como um
\textbf{curto-circuito}.\\
Pelo divisor, $G_L\to\infty$ e toda a corrente vai para o indutor:
\[
I_L=\SI{4}{\ampere};
\qquad
I_R=0.
\]
\textbf{Consequência:} um indutor em paralelo com um resistor "apaga" o resistor
em CC. Isso é usado deliberadamente para curto-circuitar componentes em regime e
deixá-los ativos apenas durante transitórios.\\
\emph{(Um capacitor em paralelo faria o oposto: em CC ele é um aberto e toda a
corrente iria ao resistor.)}}
\end{questao}

\begin{questao}[2]{}
Um divisor de corrente tem três ramos e mediu-se $\SI{4}{\ampere}$,
$\SI{2}{\ampere}$ e $\SI{1.33}{\ampere}$, com tensão comum de
$\SI{20}{\volt}$. Verifique a consistência e determine as resistências.
\resposta{Consistente; $\SI{5}{\ohm}$, $\SI{10}{\ohm}$ e $\SI{15}{\ohm}$}
\resolucao{\textbf{Resistências:} $R_k=U/I_k$:
\[
\frac{20}{4}=5;\qquad \frac{20}{2}=10;\qquad \frac{20}{1{,}33}=\SI{15}{\ohm}.
\]
\textbf{Consistência:} $\sum G=0{,}2+0{,}1+0{,}0667=\SI{0.3667}{\siemens}$, logo
$I_T=20(0{,}3667)=\SI{7.33}{\ampere}$, que confere com
$4+2+1{,}33$ \checkmark.\\
\textbf{Dois testes independentes:} a tensão comum deve ser a mesma em todos os
ramos (senão as leituras são incompatíveis) e a LKC deve fechar. Se apenas o
segundo fechasse, o erro estaria na leitura de tensão.}
\end{questao}

\begin{questao}[2]{}
Explique por que a fórmula $I_1=I_T\dfrac{R_2}{R_1+R_2}$ \textbf{não} se
generaliza para três ramos, e mostre com um contraexemplo.
\resposta{A forma correta usa condutâncias; a versão com resistências vale só
para dois ramos}
\resolucao{\textbf{Contraexemplo.} Ramos de $\SI{6}{\ohm}$, $\SI{12}{\ohm}$ e
$\SI{12}{\ohm}$ com $I_T=\SI{12}{\ampere}$.\\
\emph{Forma correta:} $I_1=12\cdot\dfrac{1/6}{1/3}=\SI{6}{\ampere}$.\\
\emph{"Generalização" ingênua} $I_1=I_T\dfrac{R_2+R_3}{R_1+R_2+R_3}$:
$12\cdot\dfrac{24}{30}=\SI{9.6}{\ampere}$ --- errado, e nem sequer a soma das
três parcelas assim calculadas daria $\SI{12}{\ampere}$.\\
\textbf{Origem da confusão:} com \emph{dois} ramos,
$\dfrac{G_1}{G_1+G_2}=\dfrac{1/R_1}{1/R_1+1/R_2}=\dfrac{R_2}{R_1+R_2}$ --- a
simplificação funciona por acaso algébrico. Com três, a manipulação equivalente
produziria $\dfrac{R_2R_3}{R_2R_3+R_1R_3+R_1R_2}$, expressão que ninguém
memoriza.\\
\textbf{Recomendação:} use sempre condutâncias. A fórmula com resistências é
uma conveniência de dois ramos, não a regra geral.}
\end{questao}

\blocoaprofundamento

\begin{questao}[3]{}
Três condutores em paralelo, de resistências $\SI{0.10}{\ohm}$,
$\SI{0.11}{\ohm}$ e $\SI{0.12}{\ohm}$, conduzem $\SI{300}{\ampere}$.
\begin{itensq}
\item Determine a corrente de cada um.
\item Determine o desvio percentual em relação à divisão ideal.
\item Comente a implicação térmica.
\end{itensq}
\resposta{$109{,}4$; $99{,}5$ e $\SI{91.2}{\ampere}$; desvios de
$+9{,}4\%$, $-0{,}5\%$ e $-8{,}8\%$}
\resolucao{(a) $\sum G=10+9{,}091+8{,}333=\SI{27.42}{\siemens}$ e
$U=300/27{,}42=\SI{10.94}{\volt}$. As correntes valem
\[
109{,}4;\qquad 99{,}5;\qquad \SI{91.2}{\ampere}.
\]
(b) Em relação a $\SI{100}{\ampere}$ ideais: $+9{,}4\%$, $-0{,}5\%$ e
$-8{,}8\%$.\\
(c) \textbf{Implicação térmica.} A dissipação vai com $I^{2}$: o condutor mais
carregado dissipa $(109{,}4/100)^{2}=1{,}20$ vezes o previsto --- $20\%$ a
mais. Ele aquece mais que os outros dois.\\
\textbf{E aí surge a pergunta decisiva:} o cobre tem coeficiente térmico
\emph{positivo}. Aquecendo, sua resistência sobe, ele devolve corrente aos
vizinhos e o sistema se \textbf{estabiliza}. Se o material tivesse coeficiente
negativo, ocorreria o contrário --- fuga térmica.\\
\textbf{Prática de instalação:} condutores em paralelo devem ter mesmo
comprimento, seção, material e trajeto. Uma diferença de $\SI{20}{\percent}$ no
comprimento produz o desequilíbrio calculado aqui, e é por isso que a norma
exige lances idênticos.}
\end{questao}

\begin{questao}[3]{}
Uma corrente de $\SI{12}{\ampere}$ deve dividir-se na razão $1{:}2{:}3$.
\begin{itensq}
\item Determine as três correntes.
\item Determine as resistências, adotando a tensão comum de
      $\SI{12}{\volt}$.
\item Determine a potência de cada ramo.
\end{itensq}
\resposta{$2$, $4$ e $\SI{6}{\ampere}$; $\SI{6}{\ohm}$, $\SI{3}{\ohm}$ e
$\SI{2}{\ohm}$; $24$, $48$ e $\SI{72}{\watt}$}
\resolucao{(a) As frações são $\frac16$, $\frac26$ e $\frac36$ de
$\SI{12}{\ampere}$: $2$, $4$ e $\SI{6}{\ampere}$.\\
(b) Como $U$ é comum, $R_k=U/I_k$:
\[
\frac{12}{2}=6;\qquad \frac{12}{4}=3;\qquad \frac{12}{6}=\SI{2}{\ohm}.
\]
As resistências ficam na razão $6{:}3{:}2$ --- \emph{inversa} de
$1{:}2{:}3$, como esperado.\\
(c) $P_k=UI_k$: $24$, $48$ e $\SI{72}{\watt}$; total $\SI{144}{\watt}=12(12)$
\checkmark.\\
\textbf{Observação de projeto:} a tensão comum foi \emph{escolhida}. Adotando
$\SI{1.2}{\volt}$ em vez de $\SI{12}{\volt}$, as resistências cairiam por dez
($0{,}6$; $0{,}3$; $\SI{0.2}{\ohm}$) e a dissipação total também --- de
$144$ para $\SI{14.4}{\watt}$. A razão de divisão é a mesma; o que muda é a
perda. Esse grau de liberdade é explorado no bloco de desafio.}
\end{questao}

\begin{questao}[3]{}
Dois ramos nominalmente iguais de $\SI{10}{\ohm}$, com tolerância de
$\pm5\%$, dividem $\SI{10}{\ampere}$.
\begin{itensq}
\item Determine o pior desequilíbrio de corrente.
\item Compare o desvio percentual da corrente com o dos resistores.
\end{itensq}
\resposta{$5{,}25$ e $\SI{4.75}{\ampere}$; desvio de $\pm5\%$}
\resolucao{(a) O pior caso ocorre com desvios opostos: $R_1=\SI{9.5}{\ohm}$ e
$R_2=\SI{10.5}{\ohm}$.
\[
I_1=10\cdot\frac{10{,}5}{20}=\SI{5.25}{\ampere};
\qquad
I_2=\SI{4.75}{\ampere}.
\]
(b) Em relação aos $\SI{5}{\ampere}$ ideais, o desvio é de
$\pm5\%$ --- \emph{igual} à tolerância dos resistores.\\
\textbf{Por que dá exatamente isso:} a sensibilidade relativa da divisão é
$\dfrac{R_1}{R_1+R_2}=0{,}5$, e o pior caso combina as duas tolerâncias:
$0{,}5\times(5\%+5\%)=5\%$.\\
\textbf{Generalização:} com ramos \emph{desiguais}, o fator muda. Em um divisor
$1{:}9$, a sensibilidade do ramo pequeno é $0{,}9$ e o desequilíbrio chega a
$9\%$ --- o ramo minoritário é sempre o mais afetado por tolerâncias.}
\end{questao}

\begin{questao}[3]{}
Um ramo de $\SI{0.5}{\ohm}$ de um divisor adquire resistência de contato de
$\SI{0.05}{\ohm}$ em suas conexões. O outro ramo, também de
$\SI{0.5}{\ohm}$, permanece íntegro, e a corrente total é
$\SI{100}{\ampere}$.
\begin{itensq}
\item Determine as novas correntes.
\item Determine o desvio e a potência adicional no contato.
\end{itensq}
\resposta{$\SI{47.6}{\ampere}$ e $\SI{52.4}{\ampere}$; desvio de
$\mp4{,}8\%$; $\SI{113}{\watt}$ no contato}
\resolucao{(a) O ramo afetado passa a $\SI{0.55}{\ohm}$:
\[
I_1=100\cdot\frac{0{,}50}{1{,}05}=\SI{47.6}{\ampere};
\qquad
I_2=100\cdot\frac{0{,}55}{1{,}05}=\SI{52.4}{\ampere}.
\]
(b) Desvio de $-4{,}8\%$ e $+4{,}8\%$ em relação aos
$\SI{50}{\ampere}$ ideais.\\
No contato: $P=0{,}05(47{,}6)^{2}=\SI{113}{\watt}$ --- concentrados em uma área
minúscula.\\
\textbf{Por que é grave.} A resistência de contato \emph{cresce} com a corrosão
e com a temperatura. Mais resistência $\to$ mais aquecimento $\to$ mais
oxidação $\to$ mais resistência: realimentação positiva que termina em falha
térmica. Curiosamente, o desvio de corrente é \emph{estabilizador} (o ramo ruim
conduz menos), mas a concentração de potência no ponto de contato domina o
processo.\\
\textbf{Conclusão:} em divisores de alta corrente, a qualidade das conexões
importa mais que a tolerância dos resistores.}
\end{questao}

\begin{questao}[3]{}
Uma corrente de $\SI{10}{\ampere}$ divide-se entre $\SI{5}{\ohm}$ e
$\SI{20}{\ohm}$.
\begin{itensq}
\item Determine as correntes e as potências.
\item Determine qual ramo aquece mais e por qual fator.
\end{itensq}
\resposta{$8$ e $\SI{2}{\ampere}$; $320$ e $\SI{80}{\watt}$ --- o menor
resistor dissipa $4\times$ mais}
\resolucao{(a)
\[
I_1=10\cdot\frac{20}{25}=\SI{8}{\ampere};
\qquad
I_2=\SI{2}{\ampere};
\]
\[
P_1=5(8)^{2}=\SI{320}{\watt};
\qquad
P_2=20(2)^{2}=\SI{80}{\watt}.
\]
(b) O ramo de $\SI{5}{\ohm}$ dissipa quatro vezes mais.\\
\textbf{Razão algébrica:} com $U$ comum, $P=U^{2}/R$, logo
$P_1/P_2=R_2/R_1=4$. A razão das potências é a \emph{inversa} da razão das
resistências --- e igual à razão das correntes ($8/2=4$).\\
\textbf{Erro comum:} supor que o resistor maior aquece mais. Isso vale em
\emph{série} ($P\propto R$), não em paralelo. Em um divisor de corrente, o
componente crítico para dimensionamento térmico é sempre o de \emph{menor}
resistência.}
\end{questao}

\begin{questao}[3]{}
Três ramos de $\SI{10}{\ohm}$, $\SI{20}{\ohm}$ e $\SI{20}{\ohm}$ têm o
primeiro aberto. Compare o efeito sob dois tipos de alimentação:
\begin{itensq}
\item fonte de corrente ideal de $\SI{12}{\ampere}$;
\item fonte de tensão ideal de $\SI{60}{\volt}$.
\end{itensq}
\resposta{(a) tensão dobra ($60\to\SI{120}{\volt}$), ramos vão de $3$ a
$\SI{6}{\ampere}$; (b) tensão e ramos inalterados, total cai de $12$ a
$\SI{6}{\ampere}$}
\resolucao{\textbf{(a) Fonte de corrente.} Antes:
$\sum G=\SI{0.2}{\siemens}$, $U=\SI{60}{\volt}$,
$I=6;3;\SI{3}{\ampere}$.\\
Depois: $\sum G=\SI{0.1}{\siemens}$, e a fonte \emph{eleva a tensão} para
manter $\SI{12}{\ampere}$:
\[
U=\frac{12}{0{,}1}=\SI{120}{\volt};
\qquad
I=6;\ \SI{6}{\ampere}.
\]
Cada ramo remanescente \textbf{dobra} de corrente, e a potência quadruplica
(de $180$ para $\SI{720}{\watt}$ em cada).\\
\textbf{(b) Fonte de tensão.} $U=\SI{60}{\volt}$ permanece.
Antes: $6;3;\SI{3}{\ampere}$, total $\SI{12}{\ampere}$. Depois:
$3;\SI{3}{\ampere}$, total $\SI{6}{\ampere}$. \textbf{Os ramos remanescentes
não sentem nada.}\\
\textbf{Síntese.} A mesma falha, no mesmo circuito, produz sobrecarga severa em
um caso e nenhum efeito no outro. O que decide é a natureza da fonte:
\begin{itemize}
\item fonte de \emph{tensão} $\Rightarrow$ ramos desacoplados, falha localizada;
\item fonte de \emph{corrente} $\Rightarrow$ ramos acoplados, falha propagada.
\end{itemize}
É por isso que instalações elétricas são alimentadas por tensão, e por que
circuitos em série de LEDs acionados por corrente exigem proteção contra
abertura.}
\end{questao}

\begin{questao}[3]{}
Um ramo fixo de $\SI{10}{\ohm}$ está em paralelo com um reostato de
$\SI{0}{\ohm}$ a $\SI{40}{\ohm}$, sob corrente total de $\SI{5}{\ampere}$.
\begin{itensq}
\item Determine a faixa de corrente no ramo fixo.
\item Determine a posição do reostato para divisão igual.
\item Comente a linearidade do controle.
\end{itensq}
\resposta{(a) de $0$ a $\SI{4}{\ampere}$; (b) $R=\SI{10}{\ohm}$}
\resolucao{(a) Com o reostato em $\SI{0}{\ohm}$, ele curto-circuita o ramo
fixo: $I_{fixo}=0$. Com $\SI{40}{\ohm}$:
\[
I_{fixo}=5\cdot\frac{40}{50}=\SI{4}{\ampere}.
\]
A faixa é $0$ a $\SI{4}{\ampere}$ --- \emph{nunca} os $\SI{5}{\ampere}$
completos, pois o reostato sempre desvia alguma corrente.\\
(b) Divisão igual exige resistências iguais: $R=\SI{10}{\ohm}$, um quarto do
curso.\\
(c) \textbf{Fortemente não linear.} A relação
$I_{fixo}=5R/(R+10)$ é hiperbólica:
\begin{center}
\begin{tabular}{ccc}
\hline
$R$ & posição do curso & $I_{fixo}$\\
\hline
$\SI{10}{\ohm}$ & $25\%$ & $\SI{2.5}{\ampere}$\\
$\SI{20}{\ohm}$ & $50\%$ & $\SI{3.33}{\ampere}$\\
$\SI{40}{\ohm}$ & $100\%$ & $\SI{4.0}{\ampere}$\\
\hline
\end{tabular}
\end{center}
Metade da faixa útil de corrente é percorrida no primeiro quarto do curso --- o
ajuste é grosseiro no início e quase inerte no fim.}
\end{questao}

\begin{questao}[3]{}
Deduza a expressão da corrente em um ramo $R_k$ quando todos os demais são
substituídos por sua resistência equivalente $R_p$, e verifique com o divisor de
$6$, $12$ e $\SI{12}{\ohm}$ sob $\SI{12}{\ampere}$.
\resposta{$I_k=I_T\dfrac{R_p}{R_k+R_p}$; para $R_k=\SI{6}{\ohm}$ e
$R_p=\SI{6}{\ohm}$, $I_k=\SI{6}{\ampere}$}
\resolucao{\textbf{Dedução.} Agrupando todos os ramos exceto o $k$-ésimo em
$R_p$, sobram \emph{dois} ramos em paralelo --- e aí a fórmula de dois ramos é
válida:
\[
I_k=I_T\frac{R_p}{R_k+R_p},
\qquad\text{com }
\frac{1}{R_p}=\sum_{j\ne k}\frac{1}{R_j}.
\]
\textbf{Verificação.} Para $R_k=\SI{6}{\ohm}$, os demais dão
$R_p=12\parallel12=\SI{6}{\ohm}$:
\[
I_k=12\cdot\frac{6}{6+6}=\SI{6}{\ampere}.\ \checkmark
\]
\textbf{Utilidade:} é a maneira correta de \emph{recuperar} a fórmula de dois
ramos em uma rede com muitos ramos --- e mostra exatamente onde a
"generalização" ingênua erra: $R_p$ é o \emph{paralelo} dos demais, não a
\emph{soma}. No exemplo, a soma daria $\SI{24}{\ohm}$ e a corrente
$\SI{9.6}{\ampere}$, valor incorreto.}
\end{questao}

\begin{questao}[3]{}
Um divisor de corrente tem um ramo contendo uma \textbf{fonte de tensão} de
$\SI{10}{\volt}$ em série com $\SI{5}{\ohm}$; o outro ramo é um resistor de
$\SI{5}{\ohm}$. A corrente total injetada é $\SI{4}{\ampere}$.
\begin{itensq}
\item Explique por que o divisor simples não se aplica.
\item Determine as correntes.
\end{itensq}
\resposta{$I_1=\SI{1}{\ampere}$ e $I_2=\SI{3}{\ampere}$}
\resolucao{(a) O divisor de corrente pressupõe ramos \textbf{puramente
resistivos}, submetidos à mesma tensão e obedecendo $I=UG$. Um ramo com fonte
segue $I=(U-E)G$ --- relação \emph{afim}, não proporcional. A repartição deixa
de depender só das condutâncias.\\
(b) Chamando $U$ a tensão comum:
\[
\underbrace{\frac{U-10}{5}}_{I_1}+\underbrace{\frac{U}{5}}_{I_2}=4
\;\Longrightarrow\;
2U-10=20
\;\Longrightarrow\;
U=\SI{15}{\volt},
\]
\[
I_1=\frac{15-10}{5}=\SI{1}{\ampere};
\qquad
I_2=\frac{15}{5}=\SI{3}{\ampere}.
\]
\textbf{Compare:} o divisor simples, com dois ramos iguais, preveria
$\SI{2}{\ampere}$ em cada. A fonte de $\SI{10}{\volt}$ \emph{se opõe} à corrente
em seu ramo e desvia $\SI{1}{\ampere}$ para o outro.\\
\textbf{Método correto:} volte à LKC com a tensão comum como incógnita --- ou
seja, análise nodal. O divisor de corrente é um atalho, e atalhos têm
hipóteses.}
\end{questao}

\blocodesafio

\begin{questao}[4]{}
\textbf{(Demonstração geral.)} Deduza a fórmula do divisor de corrente para
$n$ ramos em paralelo.
\begin{itensq}
\item Obtenha $I_k$ em função das condutâncias.
\item Mostre que a forma com resistências vale apenas para $n=2$.
\item Verifique que $\sum I_k=I_T$ identicamente.
\end{itensq}
\resposta{$I_k=I_T\dfrac{G_k}{\sum_j G_j}$}
\resolucao{(a) Todos os ramos compartilham a tensão $U$. Pela LKC,
\[
I_T=\sum_j I_j=\sum_j UG_j=U\sum_j G_j
\;\Longrightarrow\;
U=\frac{I_T}{\sum_j G_j}.
\]
Logo
\[
I_k=UG_k=I_T\frac{G_k}{\sum_j G_j}.
\]
(b) Para $n=2$, substituindo $G=1/R$:
\[
\frac{G_1}{G_1+G_2}=\frac{1/R_1}{\dfrac{R_1+R_2}{R_1R_2}}
=\frac{R_2}{R_1+R_2}.
\]
A simplificação depende de o denominador ter \emph{exatamente} dois termos.
Para $n=3$, o resultado análogo seria
\[
I_1=I_T\frac{R_2R_3}{R_2R_3+R_1R_3+R_1R_2},
\]
que não tem a forma "resistência do outro sobre a soma" --- e para $n$ geral
envolve todos os produtos de $n-1$ resistências.\\
(c)
\[
\sum_k I_k=I_T\frac{\sum_k G_k}{\sum_j G_j}=I_T.
\]
A LKC é satisfeita \textbf{identicamente}, para quaisquer condutâncias --- ela
não é um resultado do cálculo, mas uma consequência da forma da expressão.
Verificar a soma, portanto, não detecta erro nas condutâncias; ela só detecta
erro de aritmética.}
\end{questao}

\begin{questao}[4]{}
\textbf{(Projeto de balanceamento.)} Deve-se dividir $\SI{60}{\ampere}$
igualmente entre três ramos, com perda total inferior a $\SI{10}{\watt}$.
\begin{itensq}
\item Determine a resistência de cada ramo e a queda de tensão.
\item Verifique a perda e determine a margem.
\item Discuta o compromisso entre precisão da divisão e perda.
\end{itensq}
\resposta{$R\le\SI{8.33}{\milli\ohm}$; queda de $\SI{0.167}{\volt}$;
$\SI{3.33}{\watt}$ por ramo no limite}
\resolucao{(a) Divisão igual exige três resistências iguais, com
$\SI{20}{\ampere}$ cada. A perda total é
\[
P=3RI_k^{2}=3R(400)=1200R\le10
\;\Longrightarrow\;
R\le\SI{8.33}{\milli\ohm}.
\]
Adotando $R=\SI{8.33}{\milli\ohm}$, a queda é
$U=8{,}33\,\mathrm{m}\times20=\SI{0.167}{\volt}$.\\
(b) No limite, cada ramo dissipa $\SI{3.33}{\watt}$ --- exigindo resistores de
pelo menos $\SI{10}{\watt}$ para operar com folga, e provavelmente ventilados.\\
(c) \textbf{O compromisso central.} O desequilíbrio relativo entre os ramos é
governado pela \emph{razão} entre a dispersão dos resistores de balanceamento e
as resistências \emph{parasitas} (contatos, trechos de barramento, cabos):
\begin{itemize}
\item $R$ \textbf{grande}: as parasitas tornam-se desprezíveis em comparação, e
a divisão fica precisa --- mas a perda cresce com $R$.
\item $R$ \textbf{pequeno}: pouca perda, mas as parasitas passam a dominar e o
balanceamento se degrada.
\end{itemize}
Com $R=\SI{8.33}{\milli\ohm}$ e contatos de $\SI{1}{\milli\ohm}$, a incerteza
de divisão já é da ordem de $12\%$ --- o resistor de balanceamento praticamente
não cumpre sua função.\\
\textbf{Conclusão:} não se pode ter simultaneamente divisão precisa e perda
desprezível com resistores passivos. Aplicações críticas usam balanceamento
\emph{ativo} (realimentação de corrente por ramo), que rompe o compromisso.}
\end{questao}

\begin{questao}[4]{}
\textbf{(Amplificação de desequilíbrio.)} Em $n$ condutores nominalmente
iguais, um deles tem resistência $R(1-\delta)$.
\begin{itensq}
\item Deduza a fração de corrente que ele conduz.
\item Obtenha a expressão aproximada para $\delta$ pequeno e identifique o fator
      de amplificação.
\item Avalie para $n=3$ e $\delta=5\%$, $10\%$ e $20\%$.
\end{itensq}
\resposta{Desvio relativo $\approx\dfrac{n-1}{n}\,\delta$; para $n=3$:
$3{,}4\%$, $7{,}1\%$ e $15{,}4\%$}
\resolucao{(a) As condutâncias são $\dfrac{1}{R(1-\delta)}$ e
$(n-1)$ vezes $\dfrac1R$:
\[
\frac{I_1}{I_T}=\frac{\dfrac{1}{1-\delta}}{\dfrac{1}{1-\delta}+(n-1)}
=\frac{1}{1+(n-1)(1-\delta)}.
\]
(b) Para $\delta\ll1$, expandindo em torno de $\delta=0$ (onde a fração vale
$1/n$):
\[
\frac{I_1}{I_T}\approx\frac1n\left(1+\frac{n-1}{n}\,\delta\right)
\;\Longrightarrow\;
\frac{\Delta I_1}{I_1}\approx\frac{n-1}{n}\,\delta.
\]
O \textbf{fator de amplificação} é $\dfrac{n-1}{n}$: sempre menor que $1$, mas
crescente com $n$ --- tende a $1$ para muitos condutores.\\
(c) Com $n=3$, o fator é $2/3$:
\begin{center}
\begin{tabular}{ccc}
\hline
$\delta$ & exato & aproximado\\
\hline
$5\%$ & $+3{,}45\%$ & $+3{,}33\%$\\
$10\%$ & $+7{,}14\%$ & $+6{,}67\%$\\
$20\%$ & $+15{,}4\%$ & $+13{,}3\%$\\
\hline
\end{tabular}
\end{center}
\textbf{Duas leituras.} A aproximação é boa até cerca de $10\%$ e subestima o
desvio real --- portanto \emph{não} é conservadora, e deve ser usada com
cuidado em verificação de projeto.\\
E o desequilíbrio de corrente é sempre \emph{menor} que o de resistência (fator
$<1$), o que é uma boa notícia. A má notícia é que a \textbf{potência} vai com
$I^{2}$: um desvio de $7\%$ na corrente significa $15\%$ a mais de
aquecimento no condutor mais carregado.}
\end{questao}

\begin{questao}[4]{}
\textbf{(Sensibilidade.)} Para um divisor de dois ramos:
\begin{itensq}
\item Obtenha $\dfrac{\partial I_1}{\partial R_1}$.
\item Obtenha a sensibilidade relativa e mostre que $|S_1|+|S_2|=1$.
\item Interprete os casos $R_1\ll R_2$ e $R_1\gg R_2$.
\end{itensq}
\resposta{$S_1=-\dfrac{R_1}{R_1+R_2}$, $S_2=+\dfrac{R_2}{R_1+R_2}$}
\resolucao{(a) De $I_1=I_T\dfrac{R_2}{R_1+R_2}$:
\[
\frac{\partial I_1}{\partial R_1}=-\frac{I_TR_2}{(R_1+R_2)^{2}}.
\]
(b) Sensibilidade relativa:
\[
S_1=\frac{\partial I_1}{\partial R_1}\cdot\frac{R_1}{I_1}
=-\frac{I_TR_2}{(R_1+R_2)^{2}}\cdot\frac{R_1(R_1+R_2)}{I_TR_2}
=-\frac{R_1}{R_1+R_2}.
\]
Analogamente, $S_2=+\dfrac{R_2}{R_1+R_2}$, e
\[
|S_1|+|S_2|=\frac{R_1+R_2}{R_1+R_2}=1.
\]
\textbf{Mesmo padrão} já encontrado em série, paralelo e divisor de tensão ---
consequência de $I_1/I_T$ ser função homogênea de grau zero nas resistências.\\
(c) \textbf{$R_1\ll R_2$:} $S_1\to0$. O ramo pequeno leva quase toda a corrente
e é \emph{insensível} ao próprio valor --- mas $S_2\to1$, ou seja, sua corrente
depende integralmente do ramo grande. Contraintuitivo e importante.\\
\textbf{$R_1\gg R_2$:} $S_1\to-1$. O ramo grande conduz pouco, e essa pouca
corrente é inteiramente determinada pelo seu próprio valor.\\
\textbf{Regra de projeto:} para dividir corrente com precisão, aperte a
tolerância do ramo que conduz \emph{menos} --- é ele que carrega a maior
sensibilidade relativa.}
\end{questao}

\begin{questao}[4]{}
\textbf{(Divisor de baixa perda para medição.)} Deseja-se medir
$\SI{100}{\ampere}$ desviando uma fração conhecida para um instrumento, com
perda mínima.
\begin{itensq}
\item Projete um divisor de razão $1000{:}1$ com ramo principal de
      $\SI{0.1}{\milli\ohm}$.
\item Determine a corrente de medição, a queda e a perda total.
\item Compare com um resistor \emph{shunt} direto de $\SI{1}{\milli\ohm}$.
\end{itensq}
\resposta{Ramo de medição de $\SI{100}{\milli\ohm}$; $\SI{99.9}{\milli\ampere}$;
perda de $\SI{1}{\watt}$ contra $\SI{10}{\watt}$}
\resolucao{(a) A razão de correntes é a razão inversa das resistências:
\[
\frac{I_{med}}{I_{princ}}=\frac{R_{princ}}{R_{med}}=\frac{1}{1000}
\;\Longrightarrow\;
R_{med}=1000\times0{,}1\,\mathrm{m}=\SI{100}{\milli\ohm}.
\]
(b)
\[
G_{tot}=10^{4}+10=\SI{10010}{\siemens};
\qquad
I_{med}=100\cdot\frac{10}{10010}=\SI{99.9}{\milli\ampere}.
\]
\[
R_{eq}=\SI{99.9}{\micro\ohm};
\qquad
U=100(\pot{99.9}{-6})=\SI{9.99}{\milli\volt};
\qquad
P=100(\pot{9.99}{-3})=\SI{0.999}{\watt}.
\]
(c) \textbf{Shunt direto de $\SI{1}{\milli\ohm}$:}
$U=\SI{100}{\milli\volt}$ e $P=\SI{10}{\watt}$ --- \textbf{dez vezes} mais
perda.\\
\textbf{O compromisso, explicitado.} A perda vale $P=I_T^{2}R_{eq}$, e
$R_{eq}$ é essencialmente o ramo principal. Reduzi-lo reduz a perda \emph{e} a
tensão de medição --- que já é de apenas $\SI{10}{\milli\volt}$, exigindo
amplificação de baixo ruído e cuidado com termopares parasitas nas junções.\\
\textbf{Alternativa moderna:} sensores de efeito Hall ou transformadores de
corrente medem sem inserir resistência alguma no caminho principal, eliminando a
perda por completo --- ao custo de exatidão e de resposta em CC.}
\end{questao}

\begin{questao}[4]{}
\textbf{(Estabilidade térmica da divisão.)} Dois ramos idênticos dividem uma
corrente elevada. Um deles aquece mais que o outro.
\begin{itensq}
\item Analise a evolução para material de coeficiente térmico \emph{positivo}.
\item Repita para coeficiente \emph{negativo}.
\item Determine a condição de estabilidade e cite exemplos.
\end{itensq}
\resposta{$\alpha>0$: estável (realimentação negativa); $\alpha<0$: instável
(fuga térmica)}
\resolucao{(a) \textbf{Coeficiente positivo} (cobre, manganina, a maioria dos
metais). O ramo mais quente tem sua resistência \emph{aumentada}; sua
condutância cai e ele passa a conduzir \emph{menos}. Menos corrente significa
menos aquecimento --- o desvio se corrige sozinho.\\
Trata-se de \textbf{realimentação negativa}: o sistema converge para um novo
equilíbrio, com desequilíbrio limitado.\\
(b) \textbf{Coeficiente negativo} (semicondutores, termistores NTC, carbono). O
ramo mais quente reduz sua resistência, atrai \emph{mais} corrente, aquece
ainda mais --- e o ciclo se acelera. \textbf{Realimentação positiva}, ou fuga
térmica: um ramo acaba conduzindo praticamente toda a corrente e se destrói.\\
(c) \textbf{Condição de estabilidade:} $\alpha>0$ nos ramos, ou --- se
$\alpha<0$ for inevitável --- resistências de balanceamento em série, de valor
suficiente para dominar a variação térmica.\\
\textbf{Exemplos práticos.}
\begin{itemize}
\item \textbf{Condutores em paralelo:} cobre tem $\alpha\approx+0{,}004/\si{\celsius}$
--- naturalmente estáveis.
\item \textbf{Transistores bipolares em paralelo:} o coeficiente da junção é
\emph{negativo}, e por isso cada emissor recebe um resistor de balanceamento.
Sem ele, um transistor assume toda a corrente e queima.
\item \textbf{Diodos e LEDs em paralelo:} mesmo problema --- daí a exigência de
resistor individual por ramo, e não um único resistor comum.
\item \textbf{Resistores de precisão para divisão:} usa-se manganina, que
combina $\alpha$ pequeno \emph{e} positivo.
\end{itemize}}
\end{questao}

\begin{questao}[4]{}
\textbf{(Divisor com elemento ativo.)} Um dos ramos de um divisor contém uma
fonte de corrente \textbf{dependente} que injeta $g\,U$, sendo $U$ a tensão
comum e $g=\SI{0.05}{\siemens}$. O outro ramo é um resistor de
$\SI{10}{\ohm}$, e a corrente total injetada é $\SI{6}{\ampere}$.
\begin{itensq}
\item Explique por que o divisor simples falha.
\item Determine a tensão comum e a corrente de cada ramo.
\item Determine o valor crítico de $g$.
\end{itensq}
\resposta{(b) $U=\SI{120}{\volt}$; $\SI{12}{\ampere}$ no resistor e
$\SI{6}{\ampere}$ injetados pelo ramo ativo; (c) $g_c=\SI{0.1}{\siemens}$}
\resolucao{(a) O divisor pressupõe que \emph{cada} ramo obedeça $I=UG$ com
$G>0$ fixo. A fonte dependente também é proporcional a $U$, mas com sinal que
pode \emph{retirar} corrente do nó --- ela funciona como condutância
\textbf{negativa} de valor $-g$. A soma $\sum G$ deixa de ser a soma de
condutâncias físicas.\\
(b) LKC no nó, com a fonte dependente \textbf{injetando} $gU$:
\[
\frac{U}{10}=6+gU
\;\Longrightarrow\;
U(0{,}1-0{,}05)=6
\;\Longrightarrow\;
U=\frac{6}{0{,}05}=\SI{120}{\volt},
\]
\[
I_R=\frac{120}{10}=\SI{12}{\ampere};
\qquad
I_{dep}=0{,}05(120)=\SI{6}{\ampere}.
\]
\textbf{Verificação da LKC:} entram $6$ (fonte externa) mais $6$ (fonte
dependente), totalizando $\SI{12}{\ampere}$ --- exatamente o que sai pelo
resistor \checkmark.\\
\textbf{Observe o absurdo aparente:} o resistor conduz $\SI{12}{\ampere}$, o
\textbf{dobro} da corrente injetada externamente. Nenhum divisor passivo
produziria isso, pois nele vale sempre $I_k\le I_T$.\\
(c) A condutância líquida é $\dfrac1{10}-g$, que se anula em
\[
g_c=\SI{0.1}{\siemens}.
\]
Nesse ponto o nó "flutua" e não há solução finita; acima dele, a resistência
equivalente torna-se negativa e o circuito é instável.\\
\textbf{Lição:} elementos ativos podem produzir correntes de ramo
\emph{maiores} que a corrente total injetada, ou de sinal invertido --- ambos
impossíveis em rede puramente passiva. O divisor de corrente é uma ferramenta
de redes passivas.}
\end{questao}

\begin{questao}[4]{}
\textbf{(Escolha de escala.)} Um divisor deve repartir $\SI{12}{\ampere}$ na
razão $1{:}2{:}3$. A razão fixa apenas as \emph{proporções}; resta escolher a
escala absoluta das resistências.
\begin{itensq}
\item Mostre que a perda total é proporcional à escala escolhida.
\item Determine a escala para perda de $\SI{14.4}{\watt}$ e para
      $\SI{144}{\watt}$.
\item Estabeleça o critério de escolha na prática.
\end{itensq}
\resposta{$P=I_T^{2}R_{eq}$, e $R_{eq}$ escala com o fator adotado; escalas de
$0{,}6/0{,}3/\SI{0.2}{\ohm}$ e $6/3/\SI{2}{\ohm}$}
\resolucao{(a) Multiplicando todas as resistências por $\lambda$, as
\emph{frações} $G_k/\sum G$ não mudam --- a divisão é preservada. Mas
\[
R_{eq}=\lambda R_{eq,0}
\;\Longrightarrow\;
P=I_T^{2}R_{eq}=\lambda\,I_T^{2}R_{eq,0}.
\]
A perda é \textbf{diretamente proporcional} à escala. (É a mesma propriedade de
homogeneidade de grau 1 já usada para tolerâncias em associações.)\\
(b) Com $R=6/3/\SI{2}{\ohm}$: $R_{eq}=\SI{1}{\ohm}$ e
$P=144(1)=\SI{144}{\watt}$.\\
Com $R=0{,}6/0{,}3/\SI{0.2}{\ohm}$: $R_{eq}=\SI{0.1}{\ohm}$ e
$P=\SI{14.4}{\watt}$ --- dez vezes menos, com \emph{exatamente} a mesma divisão
$2{:}4{:}6$ ampères.\\
(c) \textbf{Critério.} A escala é limitada por baixo e por cima:
\begin{itemize}
\item \textbf{Limite inferior:} resistências pequenas demais tornam-se
comparáveis às parasitas (contatos, trilhas, cabos), que passam a governar a
divisão. Regra prática: mantenha $R\ge20$ vezes a maior parasita estimada.
\item \textbf{Limite superior:} a perda $\lambda I_T^{2}R_{eq,0}$ e a queda de
tensão crescem linearmente, exigindo componentes maiores e comprometendo o
rendimento.
\end{itemize}
\textbf{Escolha:} a menor escala que ainda domine as parasitas com folga. Com
contatos de $\SI{1}{\milli\ohm}$, a escala de $0{,}6/0{,}3/\SI{0.2}{\ohm}$ é
segura (o menor ramo é $200\times$ a parasita); reduzir mais dez vezes já
deixaria a divisão à mercê da qualidade das conexões.}
\end{questao}

% END BANCO COMPLEMENTAR

\gabarito[4.2 Divisor de corrente]
